%------------------------------------------------------------------------------------
%	                   DOCUMENT CONFIGURATIONS
%------------------------------------------------------------------------------------
\documentclass[12pt,a4paper]{cau_thesis}
%------------------------------------------------------------------------------------
%	                  CLASS SETTINGS
%------------------------------------------------------------------------------------
\linkboxcolor{CAU_MAIN}
\showlinkboxes{true}
% ---------------------------------------------------------------
%                     COVER PAGE SETTINGS
% ---------------------------------------------------------------
\title{Optimal Market Making in Intraday Electricity Markets via Hawkes Process Modeling}
\pdfsubject{Deep Reinforcement Learning for Intra-day Electricity Market Trading}
\pdfkeywords{Deep Reinforcement Learning, Electricity Market, Trading, Quantitative Finance}
\author{Robert Hennings}
\thesisauthor{Robert Hennings}
\stunumber{1169810}
\stuemail{stu236320@uni-kiel.de}
\privatemail{hennings.robert@web.de}
\submissiondate{XX.09.2026}
\chair{Chair of Financial Economics at the}
\institute{Institute for Quantitative Business and Economic Research}
\faculty{Faculty of Business, Economics and Social Science}
\cooperation{in cooperation with Vattenfall Energy Trading GmbH}
\location{Kiel $\bullet$ Germany}
\typedocument{Masterthesis}
\programme{Master of Science in Quantitative Finance}
\licencse{CC BY-NC-ND 4.0}{https://creativecommons.org/licenses/by-nc-nd/4.0/}
\GitHubRepo{Public GitHub Repository}{https://github.com/RobertHennings/MasterThesis}
% ---------------------------------------------------------------
%                     ASSESSMENT COMMITTEE PAGE
% ---------------------------------------------------------------
\firstsupervisorname{Your Supervisor’s Name}
\firstsupervisortitle{Title of Your Supervisor}
\firstsupervisordepartment{Department of Your Supervisor’s Department}
\secondsupervisorname{Second Reader’s Name}
\secondsupervisortitle{Title of Second Reader}
\secondsupervisordepartment{Department of Second Reader’s Department}
%------------------------------------------------------------------------------------
%	                  REFERENCES
%------------------------------------------------------------------------------------
\addbibresource{bibliography.bib}
%------------------------------------------------------------------------------------
%	                  DOCUMENT START
%------------------------------------------------------------------------------------
\begin{document}
%------------------------------------------------------------------------------------
%	                  HARD COVER TITLE PAGE
%------------------------------------------------------------------------------------
\begin{titlepage}
    \voffset-40mm
    \begin{center}
        \vspace*{0.3in} % Leicht reduziert von 0.5in
        \rule{\textwidth}{0.4pt}\\[0.2in]
        \makeatletter
        {\Huge \textbf{\@title}} \\[0.07in] % Leicht reduziert von 0.5in
        \makeatother
        \rule{\textwidth}{0.4pt}\\[0.2in]
        {\Large by} \\[0.15in] % Leicht reduziert von 0.2in
        {\Large \thesisauthor} \\[0.4in] % Leicht reduziert von 0.5in
        {\small Submitted to the} \\[0.15in]
        {\normalsize \chair} \\[0.1in]
        {\normalsize \institute} \\
        \vspace*{-1.2cm}
        \begin{center}
            {\includegraphics[width=0.30\paperwidth, height=0.32\paperheight]{logos/siegel.pdf}}
        \end{center}
        % \vspace*{-1.4cm}
        \vspace*{-1.2cm}
        {\normalsize \faculty} \\[0.15in]
        {\Large \typedocument} \\[0.15in]
        {\fontsize{14}{16}\selectfont \programme} \\[0.2in]
        {\fontsize{14}{16}\selectfont Christian-Albrechts-University of Kiel} \\[0.5in]

        \begin{tikzpicture}[remember picture, overlay]
            \node[anchor=south, yshift=6.0cm] at (current page.south) {%
                \begin{tabular}{c}
                    \rule{\textwidth}{0.4pt}                                                 \\%[0.1in]
                    {\fontsize{14}{16}\selectfont \submissiondateCover}                      \\[0.1in]
                    {\normalsize \textcopyright{} \thesisauthor, 2026. All rights reserved.} \\[0.1in]
                    {\normalsize \location}                                     \\[-0.1in]
                    \rule{\textwidth}{0.4pt}                                                 \\%[0.1in]
                \end{tabular}
            };
        \end{tikzpicture}
    \end{center}
\end{titlepage}
%------------------------------------------------------------------------------------
%	                  TITLE PAGE
%------------------------------------------------------------------------------------
% \begin{titlepage}
    \voffset-40mm
    \parbox{1.5cm}{\resizebox*{110pt}{!}{\includegraphics{logos/cau_logo.pdf}}}\hspace{310pt}
    \parbox{1.5cm}{\resizebox*{90pt}{!}{\includegraphics{logos/cau_kiel_wiso.pdf}}}\hspace{310pt}
    \parbox{1.5cm}{\resizebox*{110pt}{!}{\includegraphics{logos/vattenfall.pdf}}}\hspace{310pt}
    \begin{tikzpicture}[remember picture,overlay]
        \node[opacity=0.08, inner sep=0pt, anchor=south east]
        at ([xshift=4.2cm, yshift=1.8cm] current page.south east)
        {\includegraphics[width=0.58\paperwidth, height=0.60\paperheight]{logos/siegel_color.pdf}};
    \end{tikzpicture}
    % \begin{tikzpicture}[remember picture,overlay]
    %     \node[opacity=0.08, inner sep=0pt, anchor=south east]
    %         at ([xshift=6.2cm, yshift=-0.8cm] current page.south east)
    %         {\includegraphics[width=0.80\paperwidth, height=0.80\paperheight]{logos/siegel.pdf}};
    % \end{tikzpicture}
    \begin{center}
        \vspace*{0.5in}
        \makeatletter
        {\Huge \textbf{\@title}} \\[0.5in]
        \makeatother
        {\Large by} \\[0.2in]
        \makeatletter
        {\Large \thesisauthor} \\[0.5in]
        \makeatother
        {\small Submitted to the} \\[0.2in]
        {\normalsize \chair} \\[0.1in]
        {\normalsize \institute} \\[0.3in]
        {\normalsize \faculty} \\[0.4in]
        {\centering \cooperation\\[0.6in]}
        {\Large \typedocument} \\[0.2in]
        {\fontsize{14}{16}\selectfont \programme} \\[0.2in]
        {\fontsize{14}{16}\selectfont Christian-Albrechts-University of Kiel} \\[0.5in]
        {\fontsize{14}{16}\selectfont \submissiondateCover} \\[0.1in]
        \makeatletter
        {\normalsize \textcopyright{} {\thesisauthor}, 2026. All rights reserved.} \\[0.1in]
        \makeatother
        {\small \maybehref{\thesislicenseURL}{\thesislicenseText}}\\[0.03in]
        {\small \maybehref{\thesisgiturl}{\thesisgitlabel}}\\[0.03in]
        % \vspace{0.3cm}
        \pagebreak
        % ---------------------------------------------------------------
        %                     ASSESSMENT COMMITTEE PAGE
        % ---------------------------------------------------------------
        \begin{center}
            \textbf{First Assessor:} \\[0.2in]
            {\large \firstsupervisorname} \\[0.2in]
            {\large \firstsupervisortitle} \\[0.2in]
            {\large \firstsupervisordepartment} \\[0.5in]
            \textbf{Second Assessor:} \\[0.2in]
            {\large \secondsupervisorname} \\[0.2in]
            {\large \secondsupervisortitle} \\[0.2in]
            {\large \secondsupervisordepartment}
        \end{center}
        \begin{flushleft}
            \vspace{9.5cm}
            \makeatletter
            \textbf{Author:} \textnormal{\mdseries \thesisauthor} \\[0.1in]
            \makeatother
            \textbf{Matriculation number:} \textnormal{\mdseries \stunumber} \\[0.1in]
            \textbf{Stu-Mail:} \textnormal{\mdseries \stuemail} \\[0.1in]
            \textbf{Mail:} \textnormal{\mdseries \privatemail} \\[0.1in]
            \textbf{Submitted:} \textnormal{\mdseries \submissiondate} \\[0.1in]
        \end{flushleft}
    \end{center}
\end{titlepage}
% \begin{titlepage}
    \enlargethispage{3cm}
    \voffset-40mm
    \begin{center}
        {\normalsize Christian-Albrechts-University of Kiel} \\%[0.1in]
        {\normalsize \chair} \\%[0.1in]
        {\normalsize \institute} \\%[0.3in]
        {\normalsize \faculty} \\%[0.4in]
    \end{center}
    \vspace*{0.3in}
    \begin{center}
        \vspace*{0.9in}
        \makeatletter
        {\Large \textbf{\@title}} \\[0.5in]
        {\textbf{\typedocument}} \\[0.2in]
        {\normalsize submitted by} \\[0.2in]
        \small{\thesisauthor} \\[0.1in]
        {\fontsize{14}{16}\selectfont \submissiondateCover} \\%[0.2in]
    \end{center}
    \vspace*{0.3in}
    \begin{center}
        {\small \maybehref{\thesislicenseURL}{\thesislicenseText}}\\[0.03in]
        {\small \maybehref{\thesisgiturl}{\thesisgitlabel}}\\[0.1in]
    \end{center}
    % \vspace*{0.3in}
    \vfill
    \begin{flushleft}
        \small{Author: \thesisauthor} \\%[0.1in]
        \small{Matriculation number: \stunumber} \\%[0.1in]
        \small{Stu-Mail: \stuemail} \\%[0.1in]
        \small{Mail: \privatemail} \\%[0.1in]
        \vspace*{0.2in}
        \small{First Assessor: \firstsupervisorname} \\%[0.1in]
        \small{Second Assessor: \secondsupervisorname} \\%[0.1in]
    \end{flushleft}
\end{titlepage}
\begin{titlepage}
    \voffset-40mm
    \begin{center}
        \vspace*{0.3in} % Leicht reduziert von 0.5in
        \makeatletter
        {\Huge \textbf{\@title}} \\[0.4in] % Leicht reduziert von 0.5in
        \makeatother
        {\Large by} \\[0.15in] % Leicht reduziert von 0.2in
        {\Large \thesisauthor} \\[0.4in] % Leicht reduziert von 0.5in
        {\small Submitted to the} \\[0.15in]
        {\normalsize \chair} \\[0.1in]
        {\normalsize \institute} \\
        \vspace*{-1.2cm}
        \begin{center}
            {\includegraphics[width=0.28\paperwidth, height=0.30\paperheight]{logos/siegel.pdf}}
        \end{center}
        \vspace*{-1.4cm}
        {\normalsize \faculty} \\[0.15in]
        {\Large \typedocument} \\[0.15in]
        {\fontsize{14}{16}\selectfont \programme} \\[0.2in]
        {\fontsize{14}{16}\selectfont Christian-Albrechts-University of Kiel} \\[0.5in]

        \begin{tikzpicture}[remember picture, overlay]
            \node[anchor=south, yshift=6.0cm] at (current page.south) {%
                \begin{tabular}{c}
                    {\fontsize{14}{16}\selectfont \submissiondateCover}                      \\[0.1in]
                    {\normalsize \textcopyright{} \thesisauthor, 2026. All rights reserved.} \\[0.1in]
                    {\small \maybehref{\thesislicenseURL}{\thesislicenseText}}               \\[0.03in]
                    {\small \maybehref{\thesisgiturl}{\thesisgitlabel}}
                \end{tabular}
            };
        \end{tikzpicture}
        % \vspace{0.3cm}
        \pagebreak
        % ---------------------------------------------------------------
        %                     ASSESSMENT COMMITTEE PAGE
        % ---------------------------------------------------------------
        \begin{center}
            \textbf{First Assessor:} \\[0.2in]
            {\large \firstsupervisorname} \\[0.2in]
            {\large \firstsupervisortitle} \\[0.2in]
            {\large \firstsupervisordepartment} \\[0.5in]
            \textbf{Second Assessor:} \\[0.2in]
            {\large \secondsupervisorname} \\[0.2in]
            {\large \secondsupervisortitle} \\[0.2in]
            {\large \secondsupervisordepartment}
        \end{center}
        \vfill
        \begin{flushleft}
            \textbf{Author: \thesisauthor} \\[0.1in]
            \textbf{Matriculation number: \stunumber} \\[0.1in]
            \textbf{Stu-Mail: \stuemail} \\[0.1in]
            \textbf{Mail: \privatemail} \\[0.1in]
            \textbf{Submitted: \submissiondate} \\[0.1in]
        \end{flushleft}
    \end{center}
\end{titlepage}
% \begin{titlepage}
    \voffset-40mm
    \begin{center}
        \vspace*{0.3in} % Leicht reduziert von 0.5in
        \makeatletter
        {\Huge \textbf{\@title}} \\[0.4in] % Leicht reduziert von 0.5in
        \makeatother
        {\Large by} \\[0.15in] % Leicht reduziert von 0.2in
        {\Large \thesisauthor} \\[0.4in] % Leicht reduziert von 0.5in
        {\small Submitted to the} \\[0.15in]
        {\normalsize \chair} \\[0.1in]
        {\normalsize \institute} \\
        \vspace*{-0.0cm}
        \begin{center}
            {\includegraphics[width=0.29\paperwidth, height=0.24\paperheight]{logos/siegel_maths.png}}
        \end{center}
        % \vspace*{-0.8cm}
        {\normalsize \faculty} \\[0.15in]
        {\Large \typedocument} \\[0.15in]
        {\fontsize{14}{16}\selectfont \programme} \\[0.2in]
        {\fontsize{14}{16}\selectfont Christian-Albrechts-University of Kiel} \\[0.5in]

        \begin{tikzpicture}[remember picture, overlay]
            \node[anchor=south, yshift=6.0cm] at (current page.south) {%
                \begin{tabular}{c}
                    {\fontsize{14}{16}\selectfont \submissiondateCover}                      \\[0.1in]
                    {\normalsize \textcopyright{} \thesisauthor, 2026. All rights reserved.} \\[0.1in]
                    {\small \maybehref{\thesislicenseURL}{\thesislicenseText}}               \\[0.03in]
                    {\small \maybehref{\thesisgiturl}{\thesisgitlabel}}
                \end{tabular}
            };
        \end{tikzpicture}
        % \vspace{0.3cm}
        \pagebreak
        % ---------------------------------------------------------------
        %                     ASSESSMENT COMMITTEE PAGE
        % ---------------------------------------------------------------
        \begin{center}
            \textbf{First Assessor:} \\[0.2in]
            {\large \firstsupervisorname} \\[0.2in]
            {\large \firstsupervisortitle} \\[0.2in]
            {\large \firstsupervisordepartment} \\[0.5in]
            \textbf{Second Assessor:} \\[0.2in]
            {\large \secondsupervisorname} \\[0.2in]
            {\large \secondsupervisortitle} \\[0.2in]
            {\large \secondsupervisordepartment}
        \end{center}
        \vfill
        \begin{flushleft}
            \textbf{Author: \thesisauthor} \\[0.1in]
            \textbf{Matriculation number: \stunumber} \\[0.1in]
            \textbf{Stu-Mail: \stuemail} \\[0.1in]
            \textbf{Mail: \privatemail} \\[0.1in]
            \textbf{Submitted: \submissiondate} \\[0.1in]
        \end{flushleft}
    \end{center}
\end{titlepage}
%------------------------------------------------------------------------------------
%	                  ABSTRACT
%------------------------------------------------------------------------------------
\hypersetup{pageanchor=true}
\pagenumbering{roman}
\pagestyle{plain}
\chapter*{Abstract} % by using * we are not showing this chapter in the table of contents
DEEP Reinforcement (RL) is one of the most successful and recent
techniques in the field of Artificial Intelligence (AI). It has been
successfully applied to various fields, including robotics, gaming,
and finance. In this paper, we explore the application of Deep RL in
the context of algorithmic trading in the electricity market. We
propose a novel approach that combines Deep RL with advanced neural
network architectures to optimize trading strategies. Our method
leverages historical market data and real-time information to make
informed trading decisions. We evaluate our approach using a
comprehensive dataset from the electricity market and compare it with
traditional trading strategies. The results demonstrate that our Deep
RL-based method outperforms conventional approaches, achieving higher
returns and lower risk. This research contributes to the growing body
of literature on AI applications in finance and provides valuable
insights for practitioners in the energy trading sector.

% \underline{Keywords}: \textbf{Reinforcement Learning}, \textbf{Deep Learning},
\begin{keywords}
    \textbf{Reinforcement Learning}, \textbf{Deep Learning},
    \textbf{Electricity Market}, \textbf{Trading},
    \textbf{Algorithmic Trading}, \textbf{Energy Trading},
    \textbf{Financial Economics}, \textbf{Quantitative Finance},
    \textbf{Machine Learning}, \textbf{Artificial Intelligence},
    \textbf{Neural Networks}, \textbf{Time Series Analysis},
    \textbf{Forecasting}, \textbf{Decision Making},
    \textbf{Optimization}, \textbf{Stochastic Optimal Control},
    \textbf{Market Making}, \textbf{Market Microstructure},
\end{keywords}

\begin{classification}
    \textbf{C61}, \textbf{C63}, \textbf{G11}, \textbf{G17}, \textbf{Q41}
\end{classification}
\clearpage
%------------------------------------------------------------------------------------
%	                  ACKNOWLEDGEMENTS
%------------------------------------------------------------------------------------
\pagestyle{plain}
\chapter*{Acknowledgements} % by using * we are not showing this chapter in the table of contents
As a fairly storable commodity, data is often cited as the new oil,
therefore I would like to express my sincere gratitude to Leon von
Essen, ETD Product Designer at the EUREX, for providing me with access to
the EPEX SPOT data, without this project would not have been possible.
I am truly thankful for his support and collaboration, even though my
internship dates back almost four years. As he is, in fact, a \textit{market maker},
him supporting research devoted to this topic is a well rounded circle.
May the volume be with him.

\clearpage
% \input{chapters/00-acknowledgements}
%------------------------------------------------------------------------------------
%	                  TABLE OF CONTENTS
%------------------------------------------------------------------------------------
\hypersetup{linktoc=all}
\tableofcontents
\markboth{~}{} % Clear header for the table of contents
%------------------------------------------------------------------------------------
%	                  LIST OF ABBREVIATIONS
%------------------------------------------------------------------------------------
\chapter*{List of Abbreviations} % by using * we are not showing this chapter in the table of contents
\markboth{~}{} % Clear header for the table of contents
\addcontentsline{toc}{chapter}{List of Abbreviations}
\begin{abbreviations_table}
    % Must be ordered alphabetically
    RES & Renewable Energy Sources \\ CID & Continuous Intraday\\ IDM &
    Intraday Market\\ ICE & Intercontinental Exchange\\ EEX & European
    Energy Exchange\\ EPEX SPOT & European Power Exchange Spot Market\\
    IDA & Intraday Auction\\ LOB & Limit Order Book\\ API & Application
    Programming Interface\\ MO & Market Order\\ LO & Limit Order\\ CA &
    Cancellations\\ HJB & Hamilton-Jacobi-Bellman Equation\\ DRL & Deep
    Reinforcement Learning\\ MLE & Maximum Likelihood Estimation\\ DA &
    Day-Ahead Market\\ GWh & Gigawatt Hour\\ CET & Central European
    Time\\ TSO & Transmission System Operator\\ SDAC & Single Day-Ahead
    Coupling\\ NEMO & Nominated Electricity Market Operator\\ PPA & Power
    Purchase Agreement\\ MTU & Market Time Unit\\ SIDC & Single Intraday
    Coupling\\ SOB & Shared Orderbook\\ MWh & Megawatt Hour\\ FOK & Fill
    or Kill\\ IOC & Immediate or Cancel\\ AON & All or None\\ MM & Market
    Maker\\ ZI & Zero Intelligence\\ DMM & Designated Market Maker\\
    i.i.d & independent and identically distributed\\ HPP & Homogeneous
    Poisson Process\\ NHPP & Non-Homogeneous Poisson Process\\ MHP &
    Marked Hawkes Process\\ SDHP & State Dependent Hawkes Process \\
    SDMHP & State Dependent Marked Hawkes Process \\ NHP & Neural Hawkes
    Process \\ KKT & Karush-Kuhn-Tucker\\ LSTM & Long Short-Term Memory\\
    L\{X\} & (Price) Level X of the LOB\\ DPP & Dynamic Programming
    Principle\\ PDE & Partial Differential Equation\\ ACF &
    Autocorrelation Function \\ PACF & Partial Autocorrelation Function
    \\ PDF & Probability Density Function \\ CDF & Cumulative Probability
    Density Function \\ ECDF & Empirical Cumulative Probability Density
    Function \\ QQ & Quantile-Quantile \\ MLE & Maximum Likelihood
    Estimation \\ EEG & Erneuerbare Energien Gesetz \\ EUPHEMIA &
    Pan-European Hybrid Electricity Market Integration Algorithm \\
\end{abbreviations_table}
%------------------------------------------------------------------------------------
%	                  LIST OF SYMBOLS
%------------------------------------------------------------------------------------
\chapter*{List of Symbols} % by using * we are not showing this chapter in the table of contents
\markboth{~}{} % Clear header for the table of contents
\addcontentsline{toc}{chapter}{List of Symbols}
\vfill
\begin{symbols_table}
    % Order Book and Trading Mechanics related abbreviations
    \textit{t}    & discrete time index t                \\
    \textit{T}    & discrete terminal time index T       \\
    $\lambda(t)$  & Hawkes Process intensity             \\
    $\delta$      & order book depth                     \\
    $\Delta p$    & tick size                            \\
    $P(\delta)$   & fill-probability of order book depth \\
    \textit{V}    & order volume                         \\
    \textit{q}    & inventory level                      \\
    $\mathcal{A}$ & Set of Admissible Strategies         \\
    $\phi$        & (Running) Inventory penalty          \\
    $\lambda$     & (Conditional) Intensity              \\
    $\mathcal{F}$ & Filtration                           \\
    \textit{S}    & Midprice                             \\
    \textit{X}    & Cash-Process                         \\
    \textit{W}    & Brownian Motion                      \\
\end{symbols_table}
\vfill
%------------------------------------------------------------------------------------
%	                  LIST OF FIGURES
%------------------------------------------------------------------------------------
\chapter*{List of Figures}
\markboth{~}{}
\addcontentsline{toc}{chapter}{List of Figures}
\listoffigures
%------------------------------------------------------------------------------------
%	                  LIST OF TABLES
%------------------------------------------------------------------------------------
\chapter*{List of Tables}
\markboth{~}{}
\addcontentsline{toc}{chapter}{List of Tables}
\listoftables
%------------------------------------------------------------------------------------
%	                  LIST OF ALGORITHMS
%------------------------------------------------------------------------------------
\chapter*{List of Algorithms}
\markboth{~}{} % Clear header for the table of contents
\addcontentsline{toc}{chapter}{List of Algorithms}
% \listofalgorithms
\listofalgorithmshere
%------------------------------------------------------------------------------------
%	                  INTRODUCTION
%------------------------------------------------------------------------------------
\chapter{Introduction}\label{ch:introduction}
\pagenumbering{arabic}
% 01-introduction
The ever increasing power generation attributed to renewable energy sources (RES),
predominantly in Germany within Europe, with a current share of 57,2 \% of produced
gross-electricity by RES (2025)\footnote{See Data from \cite{ageb_bruttostromerzeugung_2026}, last accessed: 2026-05-13 19:08.},
highlights the rapid decarbonization of the energy sector.\footnote{\cite[See ][, p.2]{eurelectric2026shortterm}}
By regulatory design, RES have a feed-in priority\footnote{See \cite{eu_red_2018} for the relevant EU directive and \cite{ec_recommendation_supportschemes_2024} for the support scheme recommendation, as well as the national EEG \cite{eeg_2023_feed_in}.}, meaning that they are guaranteed to be
integrated into the electricity grid and are given priority over other sources of electricity
generation, thereby directly impacting the price formation process.

As the RES production, as just one of many influencial external
variables, is itself heavily weather dependent, this uncertainty is
transferred into the (european) electricity markets, resulting in
increased price volatility and a growing demand for balancing the
short term imbalances as well as forecasting errors in the continuous
intraday markets (CID), where the financial trading of these
positions is carried out.
%------------------------------------------------------------------------------------
\section{The financialization of the electricity market}\label{sec:financialization}
% \section{Electronic markets for electricity}\label{sec:electronic_markets}
%------------------------------------------------------------------------------------
In Europe, the main exchanges for trading electricity are the
Intercontinental Exchange (ICE), the NordPool, the European Energy
Exchange (EEX) and the European Power Exchange (EPEX SPOT). The
latter one, acting as the reference spot market, where the short term
trading of electricity takes place, what this thesis focuses on.

These exchanges provide a platform for trading electricity contracts,
including futures, options, and spot contracts. The spot trading is
organized in a CID market and three intraday auctions (IDA), where
the CID allows for trading of electricity contracts with delivery
within the next 24 hours, while the IDA is a single auction that
takes place at a specific times each day, allowing for trading of
electricity contracts with delivery on the following day. The CID
market offers 1-hour contracts, 30-Minute contracts, and the most
granular 15-Minute contracts, covering the entire 24-hour period of
the next following day. The IDA, on the other hand, focuses on 1-hour
contracts for the next day, providing a platform for market
participants to adjust their positions based on updated information
and forecasts. More details about the market structure and trading
mechanics of the EPEX SPOT will be provided in chapter
\hyperref[sec:structure_european_power_market]{2.1 The structure of
    the European Power Market}.

The trading in these contracts is facilitated through a limit order
book (LOB), where buyers and sellers submit their orders, which are
then matched based on price and time priority. The LOB provides a
transparent view of the market, allowing participants to see the
available liquidity and make informed trading decisions. This
\textit{order-driven market} structure is a key feature of the
electricity markets, as it allows for efficient price discovery and
facilitates the trading of electricity contracts in a competitive
environment, in contrast to a \textit{quote-driven market} structure,
where market makers provide liquidity by quoting buy and sell prices
for the contracts. The dynamics of the LOB are influenced by various
factors, including the arrival of limit orders and market orders,
cancellations, and modifications of existing orders. In order for a
smooth operating and liquid market, it is crucial to have a
sufficient number of market participants, including so called market
makers, who provide liquidity by continuously quoting buy and sell
prices for the contracts, making sure incoming orders can be executed
efficiently at all times. Liquidity provision is carried out by these
market makers at a rapid speed pocketing the bid-ask spread, which is
the difference between the highest price a buyer is willing to pay
and the lowest price a seller is willing to accept. As the bid-ask
spread is usually only a few cents wide, at best, earning this gap
has to be done several thousands of times a day, which is only
possible through the use of algorithmic trading strategies, which are
often based on statistical and machine learning techniques, designed
to quickly analyze huge amounts of market data, identify trading
opportunities, and execute trades at high speeds. With an average
daily order volume of 1,242,762 orders in the EPEX SPOT CID market in
2019, this leaves no real other options, as to make use of
algorithmic trading strategies, to be able to compete in this
market.\footnote{2019 figures, see \cite{epexspot_annual_report_2019}
    for more details.} This process is essential for maintaining an
efficient and liquid market, as it allows for quick execution of
trades and helps to reduce price volatility. The presence of market
makers also encourages more participants to enter the market, as they
can be assured of liquidity and competitive pricing. In recent years,
the role of market makers has become increasingly important in the
electricity markets, with many participants often using algorithmic
trading strategies, what is also evident in the share of submitted
orders to the EPEX SPOT via Application Programming Interfaces
(APIs), that is around 65\% and represents around 47\% of traded
volume.\footnote{2019 figures, see \cite{epexspot_annual_report_2019}
    for more details, as of 20-05-2026 455 Exchange Members are
    registered and 4XX APIs.}

The discrepancy between automatically and human interacted manual
trading volume is due that the automatically submitted orders are
often smaller in size, but more frequent, while the manually
submitted orders are often larger in size, but less frequent, what
later motivates modelling decisisons with respect to variable order
sizes.

This financialization highlights the economic importance of efficient
trading strategies and risk management techniques in the electricity
markets, as market participants seek to optimize their returns while
managing the risks associated with price volatility and uncertainty.
The unique characteristics of the electricity markets, such as the
non-storability of electricity and the influence of weather
conditions on supply and demand, further emphasize the need for
sophisticated modeling approaches that can capture these dynamics
effectively. In this context, the application of Hawkes processes for
optimal market making in intraday electricity markets offers a
promising avenue for developing advanced trading strategies that can
adapt to the complex and dynamic nature of these markets, what this
thesis essentially covers.

Here the main overall approach resembles around the idea that the LOB
can be treated as stochastic system, also referred to as
\textit{zero-intelligence model} but will be enhanced by including
some priors like non-constant order intensities or dependency
structures \footnote{See \cite{jain2024limit}.}, in contrast to the
often used second main approach, treating the LOB as a strategic game
between market participants, where the actions of one participant can
influence the actions of others, leading to a complex interplay of
strategies and counter-strategies, also known as \textit{perfect
    rational agent-based model}. By modeling the LOB as a stochastic
queueing system, we can also allow for flexible modelling techniques
for the involved order types, that can be submitted \textbf{market
    orders} (MO), \textbf{Limit Orders} (LO) and \textbf{cancellations},
which are the main drivers of the LOB dynamics and are explained in
more detail in \hyperref[sub_sec:order_types]{2.2.1 Order
    Types}.\footnote{See \cite{gould2013limit}.}

% The financialization of the electricity market has led to an increase in the number of financial instruments available for trading, as well as an increase in trading volumes. This has been driven by several factors

Increasing trading volumes, Increasing renewables energy production,
increasing uncertainty through renewables through weather induced
volatility flexible planning and balancing through intraday market
mechanism, cite volume increase, connection to financial markets and
LOB, efficient price discovery, high share of algorithmic trading
unique stylized facts that need to be modelled,

% Todays electricity markets are increasingly financialized, with a
% growing number of (financial) instruments and traded volumes to manage risk,
% optimize returns and most and for all balance supply and demand to guarantee
% a stable transmission system. This financialization has been driven by several factors,
% including the liberalization of electricity markets, the increasing
% volatility of electricity prices, and the growing demand for
% renewable energy sources.
%------------------------------------------------------------------------------------
\section{Main Research Objectives of the Thesis}\label{sec:main_research_objectives}
%------------------------------------------------------------------------------------
The main objective of this thesis is to explore the application of
Hawkes processes for optimal market making in intraday electricity
markets. Specifically, we aim to develop a modeling framework that
captures the dynamics of the limit order book (LOB) and the
interactions between different order types, such as limit orders,
market orders, and cancellations. By leveraging the self-exciting
nature of Hawkes processes, we seek to model the arrival of orders
and their impact on the LOB, allowing us to develop trading
strategies that can adapt to the evolving market conditions. Through
empirical analysis and numerical simulations, we aim to evaluate the
performance of our proposed approach and compare it with traditional
market making strategies, ultimately contributing to the advancement
of trading techniques in the context of intraday electricity markets.
Taking into account the often not properly considered dynamics of
Cross excitation, time varying baseline intensities and general
intensities, external excitement through updated forecasts and
external variables, we seek to develop a comprehensive modeling
framework that can effectively capture the complex dynamics of the
LOB and provide insights into optimal market making strategies in the
context of intraday electricity markets. The main research questions
that this thesis seeks to address are:
\begin{itemize}
    \item How can Hawkes processes be effectively applied to model the dynamics
          of the limit order book in intraday electricity markets?
    \item What are the key stylized facts of the limit order book in the EPEX
          SPOT CID market, and how can they be incorporated into the modeling
          framework?
    \item How does the performance of the proposed Hawkes process-based market
          making strategy compare to traditional market making strategies in
          terms of profitability and risk management?
    \item What are the implications of the findings for practitioners in the
          energy trading sector, and how can they be used to inform trading
          strategies and risk management techniques in intraday electricity
          markets?
\end{itemize}

In answering the above objectives, the main contributions are

This thesis is aimed at obtaining the optimal, dynamic quote
distances from the best bid/ask using stochastic optimal control
theory solving the Hamilton-Jacobi-Bellman (HJB) equation. This is
done, proposing different, extending modelling frameworks, all making
use of the multivariate marked Hawkes model with time-varying
baseline intensities. The bid and ask side will be modelled
separately, allowing for interactions, as well as per side, then the
three main components, affecting price building, namely Market Order
Arrivals, Limit Order Arrivals and Cancellations.By starting from
this baseline model framework with six (interaction) channels, it
will be extended to the main model by also incorporating the first
three price levels per side. Two additional extensions will be
reviewed, namely the \textit{state-dependent} and the \textit{marked}
one. In the state-dependent extension, the intensities will be
dependent on parameters like imbalance, spread, queue size,
inventory, time to delivery and time of day. In the marked extension,
order size and price impact will be used as marks, especially for MO
and LO. A possible optional extension could be the use of Deep
Reinforcement Learning (DRL), since especially in the state-dependent
extensions, we face a rich state space, which can be hard to solve
through traditional methods, but is well suited for DRL, which can be
used to learn optimal policies. For the estimation strategy, the
maximum likelihood estimation (MLE) will be the main workhorse, but
we will consider also non-parametric spline based estimation for the
excitation kernel functions.
%------------------------------------------------------------------------------------
\section{Structure of the Thesis}\label{sec:thesis_structure}
%------------------------------------------------------------------------------------
The thesis is organised into three main parts, where the first one
comprises the theoretical background, the second one the empirical
characteristics of the LOB in the EPEX SPOT CID market and the third
one the mathematical modelling and its numerical analysis. The
structure of the thesis is outlined as follows, after having laid out
the motivation for the chosen research topic along with its main
research objectives and their added new insights in
\hyperref[ch:introduction]{1 Introduction} the theoretical background
follows. In chapter \hyperref[ch:theoretical_background]{2
    Theoretical Background} we provide the necessary theoretical details
on the structure of the european power market, its organization,
functional mechanics and the key problem of market making, what this
thesis focuses on. Following in
\hyperref[ch:empirical_characteristics]{3 Empirical LOB
    Characteristics} we present the empirical characteristics, or
\textit{stylized facts}, of the LOB in the EPEX SPOT CID market,
which act as the main orientation basis and benchmark for the
modelling approach in the later chapters.\footnote{See
    \cite{jain2024limit}.} Being equipped with the needed background
knowledge about the market mechanics and typical patterns, the
central modelling aspcets are covered in
\mbox{\hyperref[ch:mathematical_modelling]{4 Mathematical
        Modelling}}. There the Hawkes process, along with its underlying
mathematical theory, is introduced and the modelling approach for
optimal market making is discussed. As the proposed modelling is put
to test in chapter \hyperref[ch:numerical_analysis]{5 Numerical
    Analysis and Simulation} being used on actual, empirical data from
the EPEX SPOT CID market, where the performance of the optimal market
making framework is evaluated, the discussion is carried out in
chapter \hyperref[ch:discussion]{6 Discussion}. Final conclusions are
drawn in chapter \hyperref[ch:conclusion]{7 Conclusion}, where the
main findings of the thesis are summarized, their implications are
discussed, and avenues for future research are suggested.

% Essentially, a LOB and all of its states, can be modelled by the two
% central inputs, that are the arrival of limit orders and market
% orders. Limit orders are orders to buy or sell a certain quantity of
% an asset at a specified price or better, while market orders are
% orders to buy or sell a certain quantity of an asset immediately at
% the best available price. The dynamics of the LOB are driven by the
% interactions between these two types of orders, as well as
% cancellations and modifications of existing orders. By analyzing the
% patterns and behaviors of limit and market order arrivals, we can
% gain insights into the underlying market dynamics and develop
% strategies for trading and risk management.
%------------------------------------------------------------------------------------
%	                  THEORETICAL BACKGROUND
%------------------------------------------------------------------------------------
\chapter{Theoretical Background}\label{ch:theoretical_background}
We begin by introducing the structure of the European Power Market,
with a particular focus on the EPEX SPOT market, which serves as the
primary context for our study. We then delve into the organization of
trading which is based on the LOB system, where the main
order-variants and the matching mechanism will be discussed. The
chapter will be closed by analysing the key problem of market making,
which is the main focus of this work.
%------------------------------------------------------------------------------------
\section{The structure of the European Power Market}\label{sec:structure_european_power_market}
%------------------------------------------------------------------------------------
The core structure of the european wholesale power market is
fourfold, consisting of the Derivatives Market, the Day-Ahead or
Forward Market, the Intraday Market, and the Balancing
Market.\footnote{The Day-Ahead and the Intraday Market together are
    most commonly referred to as \textit{Spot Markets}, \cite[see ][,
        p.136.]{bnetza_monitoringreport_energy_2025}} Their (chronological)
order can be seen in
\hyperref[fig:european_power_market_design]{Figure 2.1}, where a
special emphasis is put on the timely ordering of the different
markets and their events, as this plays a crucial role for modelling.
Their order follows the time-to-delivery of the traded contracts,
with the Derivatives Market being the furthest away from delivery,
followed by the Day-Ahead Market, the Intraday Market, and finally
the Balancing Market, which is closest to delivery.
%------------------------------------------------------------------------------------
%	                  FIGURE - START
%------------------------------------------------------------------------------------
\begin{figure}[htbp]
    \vspace*{0.5cm}
    \centering
    \blendfigure{1.0\textwidth}{8cm}{PowerMarkets_SchematicOverview.png}
    \captionfootcite{Schematic Overview of the European Power Market Design with exemplary hourly products in chronological order.}{Based on \cite{kiesel_efc19_2019}, as of 16.05.2026.}\label{fig:european_power_market_design}
\end{figure}
%------------------------------------------------------------------------------------
%	                  FIGURE - END
%------------------------------------------------------------------------------------
In the following, these components will be briefly discussed in their
chronological order, that also resembles the markets importance based
on the Trading Volume, that can be seen in
\hyperref[fig:european_power_market_volumes]{Figure 2.2}.
%------------------------------------------------------------------------------------
%	                  FIGURE - START
%------------------------------------------------------------------------------------
% This must be split apart so that we can label the figures a) and b) separately
\begin{figure}[htbp]
    \vspace*{0.5cm}
    \centering
    \blendfigure{1.0\textwidth}{8cm}{PowerMarketVolumeCombined.png}
    \footnotesize{Panel a): Annual Trading Volume in GWh of the year 2025, split into the core structural sub-markets. Panel b): Historic development of the trading volume for the two categories Derivatives Market and Spot Market, where the latter one includes the Day-Ahead and Intraday Market volumes.}
    \captionfootcite{European Power Market Trading Volume: Historic (panel b, right) and Year 2025 (panel a, left) Figures.}{Based on \cite{kiesel_efc19_2019}, as of 16.05.2026.}\label{fig:european_power_market_volumes}
\end{figure}
%------------------------------------------------------------------------------------
%	                  FIGURE - END
%------------------------------------------------------------------------------------
In the left panel of
\hyperref[fig:european_power_market_volumes]{Figure 2.2}, we can see
the annual Trading Volume in GWh of the year 2025, split into the
core structural sub-markets, where the clear dominance of the
Derivatives Market is evident, where around 91\% of traded volumes is
concentrated. In the right panel of
\hyperref[fig:european_power_market_volumes]{Figure 2.2}, we can see
the historic development of the trading volume for the two categories
Derivatives Market and Spot Market, where the latter one includes the
Day-Ahead and Intraday Market volumes. From the historic development,
the nearly 10:1 relation of the volumes is verified.
%------------------------------------------------------------------------------------
\subsection{The Derivatives Market}\label{sub_sec:derivatives_market}
%------------------------------------------------------------------------------------
The derivatives market of the EEX is the most one in terms of trading
volume, where the traded contracts are primarily settled
financially\footnote{\cite[See ][, p.
        136]{bnetza_monitoringreport_energy_2025}}, meaning that the
contracts are not necessarily linked to a physical delivery of
electricity, but rather serve as financial instruments for hedging
and speculation over the medium term.\footnote{A full list of all
    tradable instruments and their specifications can be found
    \href{https://www.eex.com/fileadmin/EEX/Downloads/Trading/Specifications/Contract_Details/eex-contract-details-xls-data_20260407.xlsx}{here}.}
The product range includes futures, options, and other derivatives
based on electricity prices, that can be traded for various delivery
periods, such as respective next 11 days and the next 2 weekends, the
current and the next 4 weeks, the current and the next 24 months, 11
quarters and 10 years. Exchange trading can be conducted from 8:00 am
to 6:00 pm (CET). In the european landscape, alongside the EEX, also
the International Commodities Exchange (ICE) and the Nasdaq OMX are
competing for trading volumes for a broader range of delivery areas.
%------------------------------------------------------------------------------------
\subsection{The Day-Ahead Market}\label{sub_sec:day_ahead_market}
%------------------------------------------------------------------------------------
In contrast to the Derivatives Market of the EEX, where the traded
contracts are primarily settled financially, a physical delivery is
mandatory on the Day-Ahead Market (also referred to as \textit{Spot
    Market}).\footnote{\cite[See ][, p.
        136]{bnetza_monitoringreport_energy_2025}} The Day-Ahead Market is a
sealed-bid, uniform price, multi-unit implicit
auction.\footnote{\cite[See ][, p.233]{balardy2022empirical}} Its
main purpose is to already lock in certain volumes for the next day
(D). Based on newly arrived forecasts, updated information and the
need to balance out the supply and demand, market participants can
adjust their positions in the Intraday Market. Market participants
can submit so-called bidding curves, stating precisely how much power
they would want to generate or consume at each price tick between the
minimum and maximum prices at each operational subperiod of the next
following day, i.e. 24 hours, 30 Minutes and 15 Minutes, resulting in
96 subperiods in total. Submitted prices at which the produces would
want to sell power is by regulatory design set to the marginal cost
of production, that dictates a natural merit-order hierarchy of the
different production technologies. RES naturally operate at the
lowest marginal costs, often assumed to be zero, therefore occupying
the lowest price levels and the left side of the supply curve, while
conventional power plants, such as coal, gas and nuclear, operate at
higher marginal costs, occupying the higher price levels and the
right side of the supply curve. These biddingcurves have to be
submitted up until the gate closure time, which is typically around
12:00 CET on the day before delivery (D-1), allowing them to react to
the latest information. Most general Market areas are further divided
into bidding zones, which are defined by the transmission system
operators (TSOs) and represent areas with similar supply and demand
conditions. The geographical area of Germany is divided into four
control areas, according to the four TSOs, which are TenneT (DE1),
50Hertz (DE2), Amprion (DE3) and TransnetBW (DE4), that are encoded
in the send orders as can be seen in the data. The price formation is
then carried out through a uniform price auction, where the market
clearing price is determined by the intersection of the aggregated
supply and demand curves, which are formed by the submitted orders
(for every granularity of the next following day, i.e. 96 15-Min
subperiods). The market clearing price is the price at which the
total quantity of electricity demanded equals the total quantity
supplied, and all transactions are executed at this price. It also
serves as a reference price for the Intraday Market and the Balancing
Market, as it reflects the market's expectations for the next day.
All power producers that have submitted orders at or below the market
clearing price will be scheduled to produce electricity, and are
awarded the market clearing price for their production. An exemplary
supply and demand curve for the market area DE-LU can be seen in
\hyperref[fig:day_ahead_curves]{Figure 2.3}, for the delivery date of
the 28th of May, 00:00-00:15. To enhance european coupling,
optimizing cross border flow and ultimately welfare, the Day-Ahead
Market is coupled across the different bidding zones (Single Day
Ahead Coupling - SDAC), meaning that the market clearing price is
determined simultaneously across all bidding zones\footnote{Exact
    pricing is carried out by the EUPHEMIA algorithm.}, taking into
account the cross-border transmission capacities and the supply and
demand conditions in each zone, which is a crucial aspect for the
price formation process and the resulting price dynamics in the
Intraday Market, as the latter one is directly linked to the
Day-Ahead Market through the trading of the same contracts.

The exchanges landscape in Europe is here more fragmented, as there
are multiple smaller nominated electricity market operators (NEMOs)
that operate Day-Ahead Markets.

Microstructurally: batch market: aggregates liquidity in a single
point in time, it aggregates liquidity at a single point in time,
eliminating intra-auction adverse selection and any latency race (cf.
Budish et al., 2015). Its uniform price is the reference for the
entire downstream complex: futures settle against it, intraday prices
are quoted (in the empirical literature) as spreads to it, and many
support schemes and PPAs index to it. Order types include simple
hourly (now quarter-hourly) step and interpolated orders, block
orders (fill-or-kill across multiple MTUs, with linked and exclusive
variants), and, on some NEMOs, more exotic constructs (loop blocks,
big blocks, scalable complex orders); a storage order type is under
public consultation with expected SDAC go-live in Q4 2027 (ENTSO-E,
2026b).
%------------------------------------------------------------------------------------
%	                  FIGURE - START
%------------------------------------------------------------------------------------
\begin{figure}[htbp]
    \vspace*{0.5cm}
    \centering
    \blendfigure{1.0\textwidth}{9cm}{DayAheadCurves.png}
    \captionfootcite{Day-Ahead Aggregated Supply and Demand Curves for the Market Area DE-LU, for the Delivery Date: 28th of May, 00:00-00.15.}{Based on \cite{kiesel_efc19_2019}, as of 27.05.2026, 13:19:50 CET/CEST.}\label{fig:day_ahead_curves}
\end{figure}
%------------------------------------------------------------------------------------
%	                  FIGURE - END
%------------------------------------------------------------------------------------

%------------------------------------------------------------------------------------
\subsection{The Intraday Market}\label{sub_sec:intraday_market}
%------------------------------------------------------------------------------------
As the physical delivery of the traded contracts is also mandatory on
the intraday market, purely financially motivated participants seek
to avoid not-netted positions heavily, as they would have to be
balanced out in the balancing market, which is the most expensive
one, as it is the last market before delivery. The main purpose of
the Intraday Market is to balance out the previously locked in
Day-Ahead positions\footnote{\cite[See ][,
        p.238]{balardy2022empirical}} that may deviate due to Power Plant
Outages, optimisation of power plant usage (generator), optimisation
of demand (costumer), renewable energy producer – changes of
forecasts\cite{kiesel_efc16_2016}. The CID Market starts on the day
before delivery (D-1) at 15:00 CET after closing of the Day-Ahead
Market with opening the order books for all 24 hours of the next
following day resulting in 24 books in parallel (with the hour 01:
00:00 - 01:00 being the first), 30 Minutes and 15 Minutes. The single
resolutions can be traded up until 5 Minutes before delivery. Due to
this nature the contract lifetime lasts from 8 hours and 55 Minutes
(for the first hour of the next following day D) up until 31 hours
and 55 Minutes (for the last hour of the next following day D).
Engaged participants therefore can trade up to 32 contracts
simultaneously, just considering the hourly
resolution.\footnote{\cite[See ][, p.234]{balardy2022empirical}} This
natural ordering of the hours over the course of the day, also groups
them inherently into ``peak' and ``off-peak' hours, stressing their
different dynamics during their contract lifetime for trading.
Simplified examples for hourly products can be seen in
\hyperref[fig:european_power_market_design]{Figure 2.1}, a detailed
table listing all hourly specificas can be found in the appendix in
\hyperref[tab:epex_parallel_durations]{Table A.1}. The main CID
market parameters are summarized below in
\hyperref[tab:epex_cid_parameters]{Table 2.1}. The pay-as-bid CID
market is accompanied by three intraday auctions (IDA), which are
single pay-as-cleared auctions, similar to the Day-Ahead market, that
take place at specific times each day:
% Idea: Make a Table out of this bullet List
\begin{itemize}
    \item \textbf{IDA 1}: At 15:00 CET on the day before delivery (D-1).
    \item \textbf{IDA 2}: At 22:00 CET on the day before delivery (D-1).
    \item \textbf{IDA 3}: At 10:00 CET on the day of delivery (D).
\end{itemize}
These IDA serve the purpose of supporting transparency and liquidity pooling, to
provide reference prices along the CID contract lifetimes. Similar as the
Day-Ahead market, the CID market is also coupled across the different bidding
zones, through the so-called Single Intraday Coupling (SIDC), which is a continuous
trading mechanism that is based on the XBID architecture, operating on
the M7 system of the EPEX SPOT. The IDA auctions are already part of this mechanism,
where the XBID architecture collects all orders in a Shared Order Book (SOB).
In a window of 20 Minutes before and after, as well as during the IDA gate-closure the continuous
XBID market is shut down.
Trading on it for a specific contract is finally shut down 30 Minutes before delivery,
still leaving trading open for participants of the same delivery-area up until 5 Minutes before delivery.
Its official chronological start is therefore at 15:20 CET at the day before delivery D-1.
The XBID related market closing times each day:
% Idea: Make a Table out of this bullet List
\begin{itemize}
    \item \textbf{Around IDA 1}: From 14:40 to 15:20 CET on the day before delivery (D-1).
    \item \textbf{Around IDA 2}: From 21:40 to 22:20 CET on the day before delivery (D-1).
    \item \textbf{Around IDA 3}: From 09:40 to 10:20 CET on the day of delivery (D).
\end{itemize}
that are also dsiplayed in the following figure \hyperref[fig:ida_acutions_xbid_closure]{Figure 2.4}:
%------------------------------------------------------------------------------------
%	                  FIGURE - START
%------------------------------------------------------------------------------------
\begin{figure}[htbp]
    \vspace*{0.5cm}
    \centering
    \blendfigure{1.0\textwidth}{8cm}{ida_acutions_xbid_closure.png}
    \captionfootcite{Schematic Overview of the IDA Auctions and XBID Closure Times.}{Based on \cite{kiesel_efc19_2019}, as of 16.05.2026.}\label{fig:ida_acutions_xbid_closure}
\end{figure}
%------------------------------------------------------------------------------------
%	                  FIGURE - END
%------------------------------------------------------------------------------------
These periods are specifically important, since this harsh reductions
in volumes can be also seen in volume plots later on.

On a european level, the EPEX SPOT and the Nord Pool exchange
concentrate the majority of the intraday trading volumes.
% Continuous double acuction on a SOB
%------------------------------------------------------------------------------------
%	                  TABLE - START
%------------------------------------------------------------------------------------
\begin{table}[htbp]
    \centering
    \scriptsize
    \begin{tabularx}{\textwidth}{@{}>{\raggedright\arraybackslash}p{0.28\textwidth} >{\raggedright\arraybackslash}X@{}}
        \toprule
        \textbf{Parameter}  & \textbf{Value}                                                                               \\
        \midrule
        Tick Size           & 0.01 EUR/MWh                                                                                 \\
        Minimum Price       & -9999.9 EUR/MWh                                                                              \\
        Maximum Price       & 9999.9 EUR/MWh                                                                               \\
        Minimum Order Size  & 0.1 MWh                                                                                      \\
        Maximum Order Size  & 1000 MWh                                                                                     \\
        Trading Hours       & 24/7 (with specific gate closure times)                                                      \\
        Trading Days        & All calendar days (including weekends and holidays)                                          \\
        Settlement Process  & Financial settlement based on the market price at delivery time                              \\
        Price Grid          & Discrete price levels determined by the tick size                                            \\
        Time Grid           & Hourly, 30-minute, and 15-minute contracts available for trading                             \\
        Gate Closure Times  & Varying based on the delivery hour, with earlier closures for earlier hours (see Figure 2.1) \\
        Order Types Allowed & Limit Orders, Market Orders, and Cancellations                                               \\
        % Add more parameters as needed
    \end{tabularx}
    \captionfootcite{Key Parameters of the EPEX SPOT CID Market.}{Based on: .}\label{tab:epex_cid_parameters}
\end{table}
%------------------------------------------------------------------------------------
%	                  TABLE - END
%------------------------------------------------------------------------------------
% Describe the key parameters that form the trading: Tick size, minimum
% order size, maximum order size, trading hours, trading days,
% settlement process, etc. Price grid, time grid, trading hours,
% trading days, settlement process, etc. For further details see
% \hyperref[sub_sec:balancing_market]{Figure 2.1}.

%------------------------------------------------------------------------------------
\subsection{The Balancing Market}\label{sub_sec:balancing_market}
%------------------------------------------------------------------------------------
As bridge to the previous subchapter: Balancing market exists as to
balance out the deviations from the locked in Intraday Market volumes

The following table will be provided in the Appendix: Test

Sequence of the market: aggregates progessively better information,
participants trade against the arrival of information, microstructure
is of each stage is disciplined by the next stages price, sequence
also a term structure of price risk
%------------------------------------------------------------------------------------
\section{Microstructure and Organisation of Intraday Power Trading}\label{sec:microstructure_and_organisation_of_power_trading}
%------------------------------------------------------------------------------------
As the main trading takes place on-exchange and is not
broker-facilitated, the core price building mechanism is represented
by the LOB.\footnote{See \cite{bnetza_monitoringreport_energy_2025},
    where only 0.2 TWh of Intraday Trading was traded on broker platforms
    vs. 102 TWh on EPEX Spot.} Most academic research is centered around
data from highly liquid markets, such as equity or forex markets,
where the bid-ask spread is usually only the tick-size as the minimum
price-increment, that is possible. Therefore these models often only
concern the best bid and aks prices, as there is most volume and a
huge market-depth, often not entirely depleted by incoming market
orders. In stark contrast, the electricity markets are characterized
by a much wider bid-ask spread, that can be several times the
tick-size, and a much lower market depth, where incoming market
orders often deplete the entire volume at the best bid or ask price,
leading to a mid-price movement. This unique characteristic of the
electricity markets necessitates a different modeling approach that
can effectively capture the dynamics of the LOB and the resulting
price movements, which is what this thesis focuses on.
%------------------------------------------------------------------------------------
\subsection{Limit Order Books and Trading Mechanics}\label{sub_sec:limit_order_books_and_trading_mechanics}
%------------------------------------------------------------------------------------
The LOB is defined on a fixed discrete grid of prices (price levels),
where the step size of each price level is called \textit{tick or
    tick-size}. This tick-size in the CID market is 0.01 EUR/MWh, which
means that the price levels are discretized in increments of 0.01
EUR/MWh. The \textit{tick-size} will be denoted by $\Delta p$. The
respective minimum and maximum volumes for submitted orders are 0.1
MWh and 1000 MWh, which means that market participants can submit
orders for quantities within this range. The order-size will be
denoted by $V$.

The trading hours are 24/7, with specific gate closure times that
vary based on the delivery hour, as illustrated in
\hyperref[fig:european_power_market_design]{Figure 2.1}. All together
these parameters are called the resolution parameters of a LOB and
are summarised in \hyperref[tab:epex_cid_parameters]{Table 2.1}.

Starting from a basic market setup, participants can submit either
Limit Orders (LO) or Market Orders (MO). Participants, sending a MO
wish to immediately buy or sell their submitted volume at the best
price currently available. Due to this unconditional specification,
MOs are also referred to as \textit{aggressive orders}, as they are
willing to accept the current market price and reduce the resting
volumes at the best bid or ask price. Across the multiple, discrete
price levels, the best bid price is the highest price at which there
is a buy order, while the best ask price is the lowest price at which
there is a sell order.

\begin{definition}[Best bid and ask prices]
    \begin{equation}
        \text{Best Bid}(t) = \max\{p_i | q_i > 0, p_i < \text{Best Ask}(t)\}
    \end{equation}
    \begin{equation}
        \text{Best Ask}(t) = \min\{p_i | q_i > 0, p_i > \text{Best Bid}(t)\}
    \end{equation}
    \label{def:best_bid_ask_prices}
\end{definition}

The natural difference between the best bid and ask prices is called
the \textit{bid-ask spread}, which represents the immediate cost of
executing a trade aggressively\footnote{\cite[See][, p.70,
    ]{cartea2015algorithmic} \cite[See ][, p.35-36]{demsetz1968cost}} and
is a key indicator of market liquidity. The bid-ask spread is defined
as the difference between the best ask price and the best bid price,
and it reflects the cost that market participants incur when
executing a trade. A narrower bid-ask spread indicates a more liquid
market, where there is a smaller cost to trading, while a wider
bid-ask spread indicates a less liquid market, where there is a
larger cost to trading.

\begin{definition}[(Quoted) Bid-ask spread]
    \begin{equation}
        \text{Bid-ask spread}_{t} = P^{\text{bid}}_{t} - P^{\text{ask}}_{t}
    \end{equation}
    \label{def:bid_ask_spread}
\end{definition}
From the full bid-ask spread, the half spread can be derived, which
represents the cost of trading for a single side of the market, either
buying or selling.
\begin{definition}[(Quoted) Half spread]\label{def:half_spread}
    \begin{equation}
        {(H_{t})_{t \geqslant 0}} = \frac{1}{2} \cdot \text{Bid-ask spread}_{t}
    \end{equation}
\end{definition}

This means that incoming MOs can potentially deplete the entire
volume at the best bid or ask price, leading to a subsequent
mid-price movement. The mid-price is defined as the arithmetic
average of the best bid and ask prices and often treated as a proxy
for the true underlying price of the asset. It assesses how close the
price of a trade is to the ´true underlying market
price´.\footnote{\cite[See][, p.71]{cartea2015algorithmic}, of course
    in some situations the true market price may lie closer to either the
    best bid or ask than the mid-price would suggest.}

\begin{definition}[Mid-price]\label{def:mid_price}
    \begin{equation}
        (S_{t})_{t \geqslant 0} = \frac{1}{2}\cdot(P^{\text{bid}}_{t} + P^{\text{ask}}_{t})
    \end{equation}
\end{definition}

The best bid and best ask can also be expressed
as:\footnote{\cite[See ][, p.80]{luckner2021hawkes}}
\begin{equation}\label{eq:best_bid_half_spread}
    \text{Best Bid}(t) = S_{t} - H_{t}
\end{equation}
\begin{equation}\label{eq:best_ask_half_spread}
    \text{Best Ask}(t) = S_{t} + H_{t}
\end{equation}

\begin{definition}[Market Order (MO)]
    A \textit{buy} (respective \textit{sell}) market order is
    \label{def:market_order}
\end{definition}

In contrast, participants sending a LO specify a price at which they
are willing to buy or sell a certain quantity of electricity, and
these orders are added to the LOB and wait for a matching
counterparty to execute the trade. LOs are also referred to as
\textit{passive orders}, as they provide liquidity to the market by
adding volume to the LOB, rather than immediately executing against
existing orders.

\begin{definition}[Limit Order (LO)]
    A \textit{buy} (respective \textit{sell}) limit order is
    \label{def:limit_order}
\end{definition}

How the incoming MO are matched with resting LO to result in a trade,
is determined by the matching mechanism of the exchange system. Often
this algorithm is based on the price-time priority principle,
according to which orders are ranked based on their price and the
time they were submitted, just as the name already suggests
\cite{cartea2015algorithmic}. Then the MO is matched with the best
available LO in the LOB, which is the LO with the highest price for a
buy MO and the LO with the lowest price for a sell MO. If there are
multiple LO at the same price level, the MO is matched with the LO
that was submitted first, following the time priority principle. This
process continues until the MO is fully executed or there are no more
matching LO available in the LOB. When the volume of the MO exceeds
the available volume at the best price level, the remaining volume of
the MO is then matched with the next best price level, and this
process continues until the entire MO is executed or there are no
more matching LO available in the LOB. This process is also referred
to as \textit{walking the book}.

If a sell-MO is executed against a buy-LO, it is said to \textit{hit
    the bid}, while if a buy-MO is executed against a sell-LO, it is said
to \textit{lift the ask}.

Besides the discussed basic order types, there are also other
modalities available:
\begin{itemize}
    \item \textbf{Iceberg Orders}: These are large orders that are split into smaller visible portions, allowing traders to hide the true size of their order from the market.
    \item \textbf{Stop Orders}: These are orders that become active only when a certain price level is reached, allowing traders to limit their losses or lock in profits.
    \item \textbf{Fill or Kill (FOK) Orders}: These are orders that must be executed immediately in their entirety or not at all, ensuring that traders either get the full quantity they want or no trade occurs.
    \item \textbf{Immediate or Cancel (IOC) Orders}: These are orders that must be executed immediately, but any portion of the order that cannot be filled is canceled, allowing traders to get as much of their order filled as possible without waiting for the entire order to be executed.
    \item \textbf{All or None (AON) Orders}: These are orders that must be executed in their entirety, but they do not have to be executed immediately, allowing traders to ensure that they get the full quantity they want without having to worry about partial fills.
    \item \textbf{Market Sweep Orders}: These are orders that are designed to quickly execute against multiple price levels in the LOB, allowing traders to quickly fill large orders without having to worry about the time it takes for the order to be executed.
\end{itemize}

Once an order has been submitted, the trader has the following three
actions available, given that the order is placed in the LOB:
\begin{itemize}
    \item \textbf{Cancel}: the order is removed from the LOB and it ceases to exist.
    \item \textbf{Modify}: the trader can modify the price or quantity of the order, and in the case of an iceberg order, the visible quantity. When the price or quantity is modified, the order loses priority in the queue.
    \item \textbf{Deactivate/Activate}: the trader can also choose to ’hibernate’ i.e. deactivate the order. In the case of hibernation the order is not visible in the LOB and can not be matched against. The reason why someone would like to hibernate their order could be, for example, in the case of losing a connection to the exchange.
\end{itemize}

Here the table of the EU Report with Percentages

The information carried by a LOB is a multivariate time series which
records the order arrival times and volumes at each price
level.\footnote{\cite[See ][, p.13]{law2015pure}} In that context, we
here explicitly stress its state based representation, for an easier
later formulation in modelling approaches. The LOB in every state,
can essentially be modelled by the two central inputs, that are the
arrival of limit orders and market orders, as these trigger the
resting LOs and result in trade executions.\footnote{See
    \cite{wu2020single}, p.1}

Here the mathematical definition of the LOB:
\begin{definition}[LOB]
    \begin{equation}
        \text{LOB}(t) = \{(p_i, q_i, t_i) | i = 1, 2, ..., N(t)\}
    \end{equation}
    \label{def:lob}
\end{definition}

where \(p_i\) is the price level of the \(i\)-th order, \(q_i\) is
the quantity of the \(i\)-th order, \(t_i\) is the time of arrival of
the \(i\)-th order, and \(N(t)\) is the total number of orders in the
LOB at time \(t\).

Here the derived optimal depths with a surrounding solid line box:

\begin{boxedequationblock}
    \begin{equation}\label{eq:optimal_depths_box}
        \begin{aligned}
            \delta^{\mathrm{SL},*}
            &= \frac{1}{\kappa^{\mathrm{BM}}}
               + \bigl(|q-1|-|q|+1\bigr)\epsilon^{\mathrm{BM}}
               + \bigl(|q-1|-|q|-1\bigr)h \\
            &\quad + \alpha(1-2q)
               - \upsilon\bigl(t,S^{\mathrm{BM}},n,h+\epsilon^{\mathrm{BM}},S^{\mathrm{BM}}\Lambda\bigr) \\
            &\quad + \upsilon\bigl(t,n,h+\epsilon^{\mathrm{BM}},S^{\mathrm{BM}}\Lambda\bigr), \\[0.5em]
            \delta^{\mathrm{BL},*}
            &= \frac{1}{\kappa^{\mathrm{SM}}}
               + \bigl(|q+1|-|q|+1\bigr)\epsilon^{\mathrm{SM}}
               + \bigl(|q+1|-|q|-1\bigr)h \\
            &\quad + \alpha(1+2q)
               - \upsilon\bigl(t,S^{\mathrm{SM}},n,h+\epsilon^{\mathrm{SM}},S^{\mathrm{SM}}\Lambda\bigr) \\
            &\quad + \upsilon\bigl(t,n,h+\epsilon^{\mathrm{SM}},S^{\mathrm{SM}}\Lambda\bigr).
        \end{aligned}
    \end{equation}
\end{boxedequationblock}

Also must reference sparse order books as here is the case because
other LOB Models are derived from highly liquid underlyings often in
quote driven marekts and not order driven markets see here:
\cite{bergault2025simulating}

%------------------------------------------------------------------------------------
\subsection{Order types and classification}\label{sub_sec:order_types}
%------------------------------------------------------------------------------------
As already described earlier, orders can be classified into two main
categories, market orders and limit orders, based on their execution
characteristics. However, within these two broad classes, the ones of
special interest for modelling purposes are those that actually move
the mid price. To this \textit{aggressive type}, also cancellations
belong, as with a cancellation, the respective order loses its
position in the price-time priority queue, leaving the next order in
line to become the new best bid or ask price.

Here the categorization into 12 types, where we focus on those types
tha actually move the mid price (altering the best bid or ask), see
also the descriptive breakdown or orders of a day and their
percentage share.

According to \cite{law2019market}, the main order types in a limit
order book (LOB) can be categorized as can be seen in the table
below:
%------------------------------------------------------------------------------------
%	                  TABLE - START
%------------------------------------------------------------------------------------
\begin{table}[htbp]
    \centering
    \resizebox{\textwidth}{!}{%
        \begin{tabular}{|l|l|c|c|c|}
            \hline
            Order Type & Order Arrival Event                    & Bid Price Effect & Ask Price Effect & Counting Process $N$ \\
            \hline
            1          & aggressive market buy                  & 0                & +                & $N_{1}(t)$           \\
            2          & aggressive market sell                 & -                & 0                & $N_{2}(t)$           \\
            3          & aggressive limit buy                   & +                & 0                & $N_{3}(t)$           \\
            4          & aggressive limit sell                  & 0                & -                & $N_{4}(t)$           \\
            5          & aggressive limit buy cancellation      & -                & 0                & $N_{5}(t)$           \\
            6          & aggressive limit sell cancellation     & 0                & +                & $N_{6}(t)$           \\
            \hline
            7          & non-aggressive market buy              & 0                & 0                & $N_{7}(t)$           \\
            8          & non-aggressive market sell             & 0                & 0                & $N_{8}(t)$           \\
            9          & non-aggressive limit buy               & 0                & 0                & $N_{9}(t)$           \\
            10         & non-aggressive limit sell              & 0                & 0                & $N_{10}(t)$          \\
            11         & non-aggressive limit buy cancellation  & 0                & 0                & $N_{11}(t)$          \\
            12         & non-aggressive limit sell cancellation & 0                & 0                & $N_{12}(t)$          \\
            \hline
        \end{tabular}}
    \vfill
    \footnotesize\textit{+/0/- indicate whether the order type causes an increase, no change, or decrease in the bid or ask price, respectively.}
    \captionfootcite{Main Order Types in a Limit Order Book.}{Based on \cite{law2019market}, as of 16.05.2026.}\label{tab:order_types}
\end{table}
%------------------------------------------------------------------------------------
%	                  TABLE - END
%------------------------------------------------------------------------------------
Explaining the top six \emph{aggressive} order types more
closely\footnote{\cite[See ][, p.15]{oxford_hawkes_2020}}:
\begin{enumerate}
    \item The placement of a limit order inside the bid-ask spread. This will
          cause the price to decrease/in- crease if it is a sell/buy limit
          order placed below/above the best ask/bid price.
    \item The cancellation of the last remaining limit order in the first queue
          on either the bid/ask side of the book. The cancellation of such a
          sell/buy limit order will cause the price to increase/decrease.
    \item The execution of the last remaining limit order at either the best
          bid/ask price. This corresponds to the placement of a sell/buy market
          order at the best bid/ask price that consumes all of the available
          liquidity at that price and has the same effect on the price as the
          cancellation of the limit order.
\end{enumerate}
To describe ONLY the top of the book, the following system can be used (ONLY top book queues)
According to \cite{law2015pure} we can write: \(N(t) = (N_{1}(t),
N_{2}(t), \ldots, N_{12}(t))\) as the multivariate point process of
the twelve order types. It further follows that, we can infer the
absolute numbers of the shaping variables: \(M_{a}(t)\) as the buy
market orders, \(M_{b}(t)\) as the sell market orders, \(S_{a}(t)\)
as the ask price and \(S_{b}(t)\) as the bid price. As the tick size,
\(\delta\), is fixed and each (minimal) price jump is one tick, the
relations hold:
\begin{equation}
    M_{a}(t) = N_{1}(t) + N_{5}(t)
\end{equation}
\begin{equation}
    M_{b}(t) = N_{2}(t) + N_{6}(t)
\end{equation}
\begin{equation}
    S_{a}(t) = S_{a}(0) + (N_{1}(t) + N_{4}(t) - N_{12}(t)) \cdot \delta
\end{equation}
\begin{equation}
    S_{b}(t) = S_{b}(0) + (N_{11}(t) - N_{2}(t) - N_{3}(t)) \cdot \delta
\end{equation}

Therefore the top book spread evolves as:
\begin{equation}
    \text{Spread}(t) = S_{a}(t) - S_{b}(t) = \text{Spread}(0) + (N_{1}(t) + N_{4}(t) - N_{12}(t)) \cdot \delta - (N_{11}(t) - N_{2}(t) - N_{3}(t)) \cdot \delta
\end{equation}

In order to avoid a situation where the spread would became less than
$\delta$, we impose that the intensities of the aggressive LO
additions are conditional on the spread based on the following
relationship, where $\mathbb{1}$ is the indicator function:
\begin{equation}
    \lambda_{i}(t) = \lambda_{i}^{'}(t) \mathbb{1}(S(t) > \delta), \quad i = 3, 4
\end{equation}

So when the spread is equal to one tick, the intensity of aggressive
LO additions drops to zero. Now, it is clear that our model is
direction and timing consistent, but since we do not keep track of
the depths of the price levels, the model is not volume consistent.

This should serve as bridge on how the different arrivals or counting
processes of these order types impact the price. Also at what times
the different sites can jump or not. We can formulate the counting
processes of all the different event types and how they affect the
different system parameters that impact the final mid-price.
%------------------------------------------------------------------------------------
\subsection{LOB Consistency}\label{sub_sec:lob_consistency}
%-------------------------------------------------------------------------------------
See here: \cite[][, p.36]{korhonen2025market}: In order for a LOB
model to be consistent with the real-world LOB dynamics it must
satisfy the following three rules:
\begin{enumerate}
    \item \textbf{Direction consistency}:
          \begin{itemize}
              \item On the arrival of a marketable buy (sell) order, only the ask (bid)
                    price can move up (down), while the bid (ask) price has to stay
                    unchanged:
              \item On the arrival of limit buy (sell) order, only the bid (ask) price
                    can move up (down) while the ask (bid) price has to stay unchanged,
                    given that the LO arrives inside the spread. Otherwise, both the bid
                    and ask prices will remain unchanged:
              \item On the cancellation of a limit buy (sell) order, only the bid (ask)
                    price can move down (up) if the canceled order is the only order at
                    the best bid (ask), while the ask (bid) price will remain unchanged.
                    Otherwise, both the bid and ask prices will remain unchanged:
          \end{itemize}
    \item \textbf{Timing consistency}: the bid and ask prices can only move at times when orders
          arrive or are cancelled:
    \item \textbf{Volume consistency}:
          \begin{itemize}
              \item On the arrival of marketable buy (sell) order whose volume is larger
                    than the available liquidity at current best ask (bid), the price
                    moves up (down) and otherwise stays unchanged:
              \item If the volume of the canceled limit buy (sell) order is equal to the
                    volume of volume available at the current best bid (ask), the bid
                    (ask) price moves down (up), and otherwise the price stays unchanged:
          \end{itemize}
\end{enumerate}
One could argue that these rules only apply when no orders can arrive at the same exact moment, but given how the matching engine of the exchange operates the orders are always handled sequentially even in the case that they arrive on the same exact nanosecond. The conditions stated above are necessary in order to model the LOB realistically, regardless of any assumptions about the arrival rates or volume distributions. However, in practice achieving a consistent LOB model is extremely difficult since we would need to keep track of the depths of all the price levels on both sides. In the next section we present the weakly consistent LOB model which only models the best bid and ask queues, and is thus only direction and timing consistent.

%------------------------------------------------------------------------------------
\subsection{Matching Mechanism}\label{sub_sec:matching_mechanism}
%------------------------------------------------------------------------------------
\section{The Market Making Problem}\label{sec:market_making_problem}
%------------------------------------------------------------------------------------
Market Making essentially consists of posting limit orders
continuously on both sides of the LOB providing liquidity to the
market, at rapid speeds, and earning the Bid-Ask spread as
compensation.\footnote{\cite[See ][, p.1; ]{jain2025impulse}
    \cite[See ][, p.9 and 20 and also the important specification of MM
        and Liquidity Privision.]{cartea2015algorithmic}}

The rational behind this way of providing
immediacy\footnote{\cite[See ][, p.]{demsetz1968cost}} can be
motivated by considering an information based perespective on market
participants. A broad classification can be made between
\textit{informed} and \textit{uninformed} traders, where the former
have access to private information that allows them to make
(suspectively) more informed trading decisions.\footnote{\cite[See
    ][, p.30]{cartea2015algorithmic}} Market Makers (MM) are typically
considered to be \textit{uninformed} traders, as they do not have
access to private information and instead rely on the public
information available in the market. However, they are still exposed
to the risk of trading with \textit{informed} traders, who may have
an informational advantage and can exploit the MM's quotes to their
benefit. This creates a situation of \textit{information asymmetry},
where the MM is at a disadvantage compared to informed traders. These
information asymmetries directly translate into the following sources
of involved risks a MM faces:
% Maybe also rather in a table than in a bullet list
\begin{itemize}
    \item \textbf{Adverse Selection Risk}: This risk arises when the MM trades with informed traders who have access to private information that allows them to make more informed trading decisions. The MM may end up buying at a higher price or selling at a lower price than the true value of the asset, resulting in potential losses.
    \item \textbf{Inventory Risk}: This risk arises from holding an inventory of assets that may fluctuate in value due to market conditions. The MM may be forced to sell at a loss if the market moves against their position, resulting in potential losses.
    \item \textbf{Latency Risk}: This risk arises from delays in receiving and processing market data, which can result in the MM's quotes being stale and not reflecting the current market conditions. This can lead to potential losses if the MM is unable to react quickly enough to changing market conditions.
\end{itemize}

This perspective is also justified by looking at the different
participant types at the EPEX SPOT, where there are not only utilitiy
companies but also financial institutions, which are typically
considered to be more informed traders, as they have access to more
sophisticated trading strategies.

Now balancing these types of risks is the core problem of market
making, as the MM must continuously adjust their quotes to manage
their inventory and minimize their exposure to adverse selection and
latency risks, while still providing liquidity to the market. The
main appoaches on how to come up with optimal strategies for market
making can be depicted in the following table:

In this thesis, the considered MM uses LO and MO for maximising
expected terminal wealth (cash + inventory) under penalisations for
the running inventory as well for terminal inventory, by finding the
set of optimal controls u $\in \mathbb{U}$, where $\mathbb{U}$ is the
set of admissible controls consisting of a sequence of decision
variables.\footnote{\cite[See ][, p.]{luckner2021hawkes}}

%------------------------------------------------------------------------------------
%	                  TABLE - START
%------------------------------------------------------------------------------------
\begin{table}[htbp]
    \centering
    \scriptsize % Slightly smaller text size to perfectly fit high-density microstructure data
    \begin{tabularx}{\textwidth}{l >{\raggedright\arraybackslash}X >{\raggedright\arraybackslash}X >{\raggedright\arraybackslash}X >{\raggedright\arraybackslash}X}
        \toprule
        \textbf{School of Thought}                                                                                                                                                         & \textbf{Core Philosophy \& Core Assumption} & \textbf{Key Methodologies \& Mathematical Tools} & \textbf{Seminal Frameworks / Examples} & \textbf{Primary Limitation} \\
        \midrule

        \textbf{1) Economics-Based}                                                                                                                                                        &
        \textbf{Perfect Rationality}: Agents maximize utility under structural constraints. Focuses on \textit{why} liquidity is priced based on risk aversion and asymmetric information. &
        $\bullet$ Dynamic programming \newline $\bullet$ Game theory (Nash / Stackelberg) \newline $\bullet$ Expected utility functions                                                    &
        $\bullet$ \textbf{Inventory Models}: Ho-Stoll, Avellaneda-Stoikov \newline $\bullet$ \textbf{Information Models}: Glosten-Milgrom, Kyle's lambda                                   &
        Highly stylized; yields closed-form solutions that fail to capture real-world, high-frequency order book dynamics.                                                                                                                                                                                                                                         \\
        \midrule

        \textbf{2) Physics-Based}                                                                                                                                                          &
        \textbf{Zero Intelligence (ZI)}: Complex macro market allocations emerge from the statistical mechanics of non-strategic, random order flows.                                      &
        $\bullet$ Stochastic processes \newline $\bullet$ Poisson arrival rates \newline $\bullet$ Master equations / Particle physics analogies                                           &
        $\bullet$ \textbf{ZI-LOB Models}: Anufriev-Panchenko, Maslov model \newline $\bullet$ \textbf{Statistical Mechanics}: Santa Fe Institute models (Farmer et al.)                    &
        Completely ignores strategic behavior; agents do not adapt to adverse selection or toxic order flow.                                                                                                                                                                                                                                                       \\
        \midrule

        \textbf{3) Mathematically-Based}                                                                                                                                                   &
        \textbf{State Dependence}: Focuses strictly on data-driven, conditional state transitions of the actual limit order book queue.                                                    &
        $\bullet$ Markov chains / Jump processes \newline $\bullet$ Point processes (Hawkes processes) \newline $\bullet$ Queueing theory                                                  &
        $\bullet$ \textbf{Queue-Reactive Model} (Huang, Lehalle, Rosenbaum) \newline $\bullet$ \textbf{Cont-Stoikov-Talreja} Markovian framework                                           &
        High dimensionality; suffers from the ``curse of dimensionality'' when scaling up the number of tracked book levels.                                                                                                                                                                                                                                       \\
        \bottomrule
    \end{tabularx}
    \captionfootcite{Frameworks of Market Making Microstructure}{Based on: Own Concept.}\label{tab:market_making_frameworks}
\end{table}
%------------------------------------------------------------------------------------
%	                  TABLE - END
%------------------------------------------------------------------------------------
For a more detailed overview of the main schools of thought and
connected concrete concepts, see \cite{jain2024limitreview}.

The resulting concepts and models based off of these theoretical
strands often lack real world applicability, as they are based on
highly stylized assumptions or just miss important real world
caveats.

Here stole everything from: \cite{law2015pure} page 5: In this
section, we highlight some issues of existing market-making models in
the con- text of high-frequency trading. 1. In modern financial
exchanges, prices are only allowed on a predefined fixed grid called
price ticks. As a result, price is a pure-jump process and it has two
dimensions, 5 namely times and magnitudes of the jumps. Diffusion can
only approximate the magnitudes of the jumps but cannot describe the
properties related to timing of the jumps such as jump clustering. 2.
The common assumption of Poisson order arrivals is often rejected in
empirical lit- erature as order arrivals depict strong
self-excitation behavior [18–20]. 3. All models assume that price and
order arrivals are independent, which is far from the truth by
realizing that price rises with large buy market order and falls with
large sell market order. Because of adverse selection [1, 2], the
absence of this crucial dependency structure will generate large
phantom profit for the market maker and cause the average profit of
the market-making strategy to be overstated. 4. Since almost all
exchanges nowadays use the price-time2 priority, changing price or
quantity of limit orders means loss of priority. Nonetheless,
existing models all use regular control to continuously adjust the
quotes without any penalty. 5. A critical component of many existing
models is the demand/supply rate function. For example, in [8], it is
in the form while in [9], it can be expressed as . Yet the parameters
are hard to estimate since when the limit order is more than one tick
from the best quote, the execution probability is minuscule (e.g.
less than 36. For the sake of simplicity, existing models assume all
orders are of the same size. Such an assumption will mask the risk of
overtrading of the market maker. For instance, to continuously
maintain priority in the queue, the market maker may post more limit
orders in the order book than his risk tolerance. However, the
arrival of one giant market order may raise his inventory to an
unacceptable level, which can potentially lead to bankruptcy. Such
kind of risk cannot be modeled with all orders having the same size.% for E-mini S&P future [21]). In addition, the quoted price is not continuous but only allowed in a fixed grid of price ticks.
%------------------------------------------------------------------------------------
%	                  FIGURE - START
%------------------------------------------------------------------------------------
\begin{figure}[htbp]
    \vspace*{0.5cm}
    \centering
    \blendfigure{1.0\textwidth}{7cm}{EPEX_member_category.png}
    \captionfootcite{EPEX SPOT exchange member categories.}{Based on \cite{kiesel_efc19_2019}, as of 31.07.2026, 17:11:50 CET/CEST.}\label{fig:epex_member_category}
\end{figure}
%------------------------------------------------------------------------------------
%	                  FIGURE - END
%------------------------------------------------------------------------------------
In contrast to other financial markets, the trading on the EPEX SPOT
does not involve designated market maker (DMM) programs, where
specific participants are incentivized to provide liquidity to the
market.\footnote{\cite[See ][, p.230]{balardy2022empirical}}

In general, the main theoretical approaches to tackle this problem
can be categorised into the following streams:
% Here maybe even a schematic figure
\begin{enumerate}
    \item Physics based zero intelligence models
    \item Economics based rational agent models
    \item Math based stochastic optimal control models
\end{enumerate}

A good overview with a more granular perspective can be found in
\cite{jain2024limitreview} where great emphasis is placed on the
modelling approaches.

The classical problem formulation based on Ho \& Stoll (1981) is to
choose the optimal depth size for posting LO into the book (See
\cite[See ][, p.29]{cartea2015algorithmic})

Also mention the advantage of potentially getting a fee rebate from
the exchange in contrast to the taker fee that liquidity takers have
to pay for their orders

Key problem: Balancing the optimal limit order posting depth and the
fill probabilities Latency risk as the EPEX Spot is not offering co
location services, so the MM is exposed to latency risk, as they are
not able to react as fast as other market participants that are
co-located with the exchange servers. This can lead to a situation
where the MM's quotes are stale and do not reflect the current market
conditions, resulting in potential losses.

Use stochastic optimal control to tackle the classical problem of
solving for the optimal depth as distance from the midprice for
posting limit orders, also core problem: how to get the fill
probabilities or approximate them
%------------------------------------------------------------------------------------
%	                  EMPIRICAL LOB CHARACTERISTICS
%------------------------------------------------------------------------------------
\chapter{Empirical Microstructure: LOB Characteristics and Market Quality}\label{ch:empirical_characteristics}
%------------------------------------------------------------------------------------
In the following chapter, so called \textit{quality-parameters} of
the LOB will be analyzed, which are the key empirical characteristics
of the LOB. These will be used to first motivate, then calibrate and
finally evaluate the mathematical model for the derivation of the
optimal market making strategy. These \textit{quality-parameters}
include the following aspects, which we defined mathematically
earlier in section
\hyperref[sec:microstructure_and_organisation_of_power_trading]{2.2
    Microstructure and Organisation of Intraday Power
    Trading}:\footnote{\cite[See][, p.61.]{cartea2015algorithmic}}:
\begin{itemize}
    \item Interarrival times % although not directly a quality parameter and not really math. derived in the bespoke section
    \item Volumes and Volatility
    \item Spread measures \footnote{\cite[See ][, p.229]{balardy2022empirical}}
          and the Half Spread
    \item Market Depths and Trade sizes
    \item Market Impact
    \item Message activity and the relationship between cancellations,
          executions and distance from the midprice
\end{itemize}
As these components are important drivers of overall \emph{liquidity}, they are key
in forming the shape of the LOB, what will be important for later modelling considerations in
chapter \hyperref[ch:mathematical_modelling]{Mathematical Modelling}.
Before we start to analyse the empirical characteristics of the LOB, we
will first provide a brief introduction of the used data.
%------------------------------------------------------------------------------------
\section{Data and Data Processing}\label{sec:data}
%------------------------------------------------------------------------------------
In order to collect all L3 market-by-order events for the delivery
area germany, we collect all single order books from the EPEX SPOT
for the TSOs: 50Hz, Amprion, Tennet and TransnetBW for the time range
of the full quarter of Q2 2026. The single observations are
timestamped on millisecond level, which is the highest possible
granularity that can be achieved with the EPEX SPOT L3 data. A more
detailed data description can be found in the \hyperref[anhang2]{Data
    Appendix} under \hyperref[sec:lob_data_description]{LOB Data
    Description}. The data contains all limit order additions and
cancellations as well as market orders for each one-hour product. Due
to the XBID mechanism, the data also includes orders from other
delivery zones, but the origin of the orders is not included in the
data. In the following table, a first descriptive breakdown with
emphasis on the average events per day is provided, that serves as
first basis for the more detailed analysis of the intraday activity.
%------------------------------------------------------------------------------------
%	                  TABLE - START
%------------------------------------------------------------------------------------
\begin{table}[htbp]
    \centering
    \begin{tabular}{|c|c|c|c|}
        \hline
        Hour & Year & Number of (Trading) Days & Average Events per Day \\
        \hline
        0    & 2019 & 365                      & X                      \\
        1    & 2019 & 365                      & X                      \\
        ...  & ...  & ...                      & ...                    \\
        23   & 2019 & 365                      & X                      \\
        \hline
    \end{tabular}
    \caption{Data Overview: Hourly Event Counts and Averages.}\label{tab:data_overview}
\end{table}
%------------------------------------------------------------------------------------
%	                  TABLE - END
%------------------------------------------------------------------------------------

%------------------------------------------------------------------------------------
\section{Intraday activity patterns}\label{sec:intraday_activity_patterns}
%------------------------------------------------------------------------------------
Here just follow the Cartea Book for all the setcions, investigate
data and also cite the literature for these findings here for
confirmation, Lead Thesis did it as well like that
%------------------------------------------------------------------------------------
\subsection{Interarrival Times}\label{sub_sec:interarrival_times}
%------------------------------------------------------------------------------------
As already mentioned in XX, when applying a more granular analysis,
the modelling of prices requires a \emph{discrete}, rather than
\emph{continuous} viewpoint.\footnote{By interarrival times, we refer
    to the interval times between movements at either the bid or the ask,
    \cite[see ][, p.48]{cartea2015algorithmic}.}

In the below figure we can see, that the interarrival times seem to
be quite good approximated by a power law distribution with very
heavy tails.\footnote{\cite[See ][, p.48]{cartea2015algorithmic}}
%------------------------------------------------------------------------------------
%	                  FIGURE - START
%------------------------------------------------------------------------------------
\begin{figure}[htbp]
    \vspace*{0.5cm}
    \centering
    \blendfigure{1.0\textwidth}{9cm}{interarrival_times.png}
    \captionfootcite{Interarrival Times.}{Based on \cite[See ][, p.49]{cartea2015algorithmic}.}\label{fig:interarrival_times}
\end{figure}
%------------------------------------------------------------------------------------
%	                  FIGURE - END
%------------------------------------------------------------------------------------
In the below figure, we can see that changes to the bid- and ask are
not independent, but have a strong autocorrelated component,
suggesting that they appear in \emph{clusters}.\footnote{\cite[See
    ][, p.49]{cartea2015algorithmic}} Rapid changes in the bid/ask are
followed by further rapid changes, while a relatively long calm
period is also similarly followed by another calm
period.\footnote{\cite[See ][, p.49]{cartea2015algorithmic}}
%------------------------------------------------------------------------------------
%	                  FIGURE - START
%------------------------------------------------------------------------------------
\begin{figure}[htbp]
    \vspace*{0.5cm}
    \centering
    \blendfigure{1.0\textwidth}{9cm}{acf_interarrival_times.png}
    \captionfootcite{Interarrival Times ACF.}{Based on \cite[See ][, p.49]{cartea2015algorithmic}.}\label{fig:acf_interarrival_times}
\end{figure}
%------------------------------------------------------------------------------------
%	                  FIGURE - END
%------------------------------------------------------------------------------------
From this we can conclude that the interarrival times are not
independent and identically distributed (i.i.d.), but rather exhibit
a strong temporal dependence structure. This has important
implications for modeling the dynamics of the limit order book, as it
suggests that past events can provide valuable information about
future market activity.
% Result for modelling: Interarrivals are correlated and dependent -> History is relevant, Hawkes Process that uses history of not independent events unlike the Poisson Process
% Result for modelling: Interarrivals seem to follow a power law distribution with heavy tails, maybe as kernel function for the Hawkes Process (?)
%------------------------------------------------------------------------------------
\subsection{Volumes and Volatilty}\label{sub_sec:volumes_and_volatility}
%------------------------------------------------------------------------------------
The volume pattern over a trading session or lifetime of a contract
for the intraday power markets is quite distinctive from other
financial markets, that typically show a U-shape over the course of a
normal trading day. Starting off with high volumes near market
opening, a slight volume lull in the middle of the day, and a peak in
volume near market close.\footnote{By volume we refer to the trading
    volume measured in MWh.} This shape generally reflects the quote
updates based on newly arrived information. In contrast, the intraday
power markets show a different pattern, with a peak in volume near
gate closure, but also a significant increase in volume during the
last hour of trading.\footnote{With activity, also the volume keeps
    up, as 80\% of total volume is traded within the last three hours of
    a trading period \cite[See ][, p.238]{balardy2022empirical}.}This is
due to the fact that market participants have more information about
the supply and demand conditions for the delivery period as it
approaches, and they adjust their positions accordingly
(\textit{trading against information arrival}). It can also be seen
clearly that the volume picks up at around 6 to 7 hours prior to gate
closure.\footnote{This is usually the time when the morning shift
    starts for traders in Europe, as the market is open 24/7, the trading
    is done in three shifts: morning, evening and night shifts or
    alternatively, the night shift might be working from a country with a
    more suitable time-zone like Singapore}. Also algorithmns may try to
close out any remaining residual volumes at the end of the trading
sessions, what possibly leads to an increase in volumes towards gate
closure in combination with execution strategies that try to postpone
non-urgent transactions to times when spreads are narrow and more
preferable, what is also the case near gate
closure.\footnote{\cite[See ][, p.66]{cartea2015algorithmic}} This
would match the observed tightening of the spread that can be seen in
\hyperref[sub_sec:spreads]{3.2.3 Spreads}.
%------------------------------------------------------------------------------------
%	                  FIGURE - START
%------------------------------------------------------------------------------------
\begin{figure}[htbp]
    \vspace*{0.5cm}
    \centering
    \blendfigure{1.0\textwidth}{9cm}{volume_distribution_heatmap.png}
    \captionfootcite{Volume distribution for hour 14 delivery.}{Based on \cite[See ][, p.63]{cartea2015algorithmic}.}\label{fig:volume_distribution_heatmap}
\end{figure}
%------------------------------------------------------------------------------------
%	                  FIGURE - END
%------------------------------------------------------------------------------------
Other distinctive features of the volume pattern are the three spikes
in volume that occur shortly after the IDAs. During the IDAs and the
final hour of trading, XBID is closed, meaning that the LOB only
contains orders with delivery at one of the four German delivery
zones.\footnote{The schematic timeline can be followed in
    \hyperref[fig:ida_acutions_xbid_closure]{Figure 2.4}.} The effect of
this can be seen in the figure as the executed volume drops visibly
during the IDAs and spikes right after when the LOBs are joined
together and overlapping orders are executed. The Volume also seems
to exhibit a susbtantial variability, as can be seen in the below
figure.
%------------------------------------------------------------------------------------
%	                  FIGURE - START
%------------------------------------------------------------------------------------
\begin{figure}[htbp]
    \vspace*{0.5cm}
    \centering
    \blendfigure{1.0\textwidth}{9cm}{intraday_volume.png}
    \captionfootcite{Intraday Volume.}{Based on \cite[See ][, p.63]{cartea2015algorithmic}.}\label{fig:intraday_volume}
\end{figure}
%------------------------------------------------------------------------------------
%	                  FIGURE - END
%------------------------------------------------------------------------------------
To focus on the time-to-delivery dependency of the traded volumes, we
see in the below functional approach heat maps (left panels) of daily
volume for the year 2013 for AAPL and MENT. The heat-maps are
generated by first bucketing traded volume into five-minute windows
throughout the day, for every day of the year. Then, for each
five-minute bucket, we compute the distribution of volume. The
heat-map is a visualisation of the collection of these distributions
for each five-minute bucket all at once. In the figures we use
coloured lines to represent the first and third quartiles, as well as
the median.\footnote{\cite[See][, p.65.]{cartea2015algorithmic}} A
clear dependency of the volume on the time-to-delivery can be seen,
as descibed earlier but here with the advantage of being able to spot
outliers (right panels).
%------------------------------------------------------------------------------------
%	                  FIGURE - START
%------------------------------------------------------------------------------------
\begin{figure}[htbp]
    \vspace*{0.5cm}
    \centering
    \blendfigure{1.0\textwidth}{9cm}{intraday_volume_time_of_day.png}
    \captionfootcite{Intraday Volume as function of the time of day.}{Based on \cite[See ][, p.65]{cartea2015algorithmic}.}\label{fig:intraday_volume_time_of_day}
\end{figure}
%------------------------------------------------------------------------------------
%	                  FIGURE - END
%------------------------------------------------------------------------------------
Volume itself seems to have significant time-series persistance and
positive and significant `correlation' with volatility.\footnote{By
    volatility we refer to the standard deviation of the price changes
    measured in €/MWh.} Intraday volume has very large
peaks.\footnote{\cite[See ][, p.64]{cartea2015algorithmic}} This is
evident by further following the time to delivery timescale but
graphing the volatility in the below figure.
%------------------------------------------------------------------------------------
%	                  FIGURE - START
%------------------------------------------------------------------------------------
\begin{figure}[htbp]
    \vspace*{0.5cm}
    \centering
    \blendfigure{1.0\textwidth}{9cm}{intraday_volatility.png}
    \captionfootcite{Intraday Volatility.}{Based on \cite[See ][, p.63]{cartea2015algorithmic}.}\label{fig:intraday_volatility}
\end{figure}
%------------------------------------------------------------------------------------
%	                  FIGURE - END
%------------------------------------------------------------------------------------

%------------------------------------------------------------------------------------
%	                  FIGURE - START
%------------------------------------------------------------------------------------
\begin{figure}[htbp]
    \vspace*{0.5cm}
    \centering
    \blendfigure{1.0\textwidth}{9cm}{price_pattern_cent_bounds.png}
    \captionfootcite{Price pattern with cent-level bounds.}{Based on \cite[See ][, p.69]{cartea2015algorithmic}.}\label{fig:price_pattern_cent_bounds}
\end{figure}
%------------------------------------------------------------------------------------
%	                  FIGURE - END
%------------------------------------------------------------------------------------
In the above figure we can see the large accumulation of transactions
whose prices end in exact dollar-values prices, with the most
pronounced spike at the round 50 cent value, with slightly less
pronounced spikes at the round 10th of cents
values.\footnote{\cite[See ][, p.64]{cartea2015algorithmic}} As
incoming aggressive MO just hit standing LO in the LOB, is is
ultimately the MMs who deciide to post their orders at these round
values for providing liquidity.\footnote{\cite[See ][,
        p.64]{cartea2015algorithmic}}

Further going into more details and dividing the tradable hourly
products into peak and off-peak products and also analysing potential
weekend vs. during the week effects, we can see that the volume
pattern varies for these subtypes. Samuleson Effect here From
previous intraday pattern: Baseline intensity has to be dependent on
the remaining contract lifetime, therefore being non-stationary

Percentage table for the different order types to stress that onyl
few types are relevant and also focus on only the best 3 price levels
because mostly only 1 tick moved (?).
% Result for modelling: Volume is not independent but depends on own history
% Result for modelling: Volume depends on the time-to-delivery: Add that as state variable
% Result for modelling: Volume and Volatility are positively correlated
% Result for modelling: Volume and Bid-Ask Spread are negatively correlated, Bid-ask spread also as possible state variable
% Result for modelling: Majority of volumes is traded in last three hours or even just the last hours
%------------------------------------------------------------------------------------
\subsection{Spreads}\label{sub_sec:spreads}
%------------------------------------------------------------------------------------
Spreads measure the immediate cost of executing a trade
aggressively.\footnote{\cite[See ][, p.70]{cartea2015algorithmic}} In
comparison to the later analysed \emph{market impact}, that centers
more around larger trades, spreads measure the execution costs of
smaller trades in relation to the current `market
price'.\footnote{\cite[See ][, p.71]{cartea2015algorithmic}}
%------------------------------------------------------------------------------------
%	                  FIGURE - START
%------------------------------------------------------------------------------------
\begin{figure}[htbp]
    \vspace*{0.5cm}
    \centering
    \blendfigure{1.0\textwidth}{9cm}{bid_ask_spread.png}
    \captionfootcite{Bid-ask spread distribution as a function of time to delivery.}{Based on \cite[See ][, p.51]{cartea2015algorithmic}.}\label{fig:bid_ask_spread}
\end{figure}
%------------------------------------------------------------------------------------
%	                  FIGURE - END
%------------------------------------------------------------------------------------
As can be seen in the figure above, the bid-ask spread over a trading
session seems to follow an L-shape pattern, in line with the findings
of: XX. The spread decreases with time as there is more activity
closer to the delivery. This is consistent with the findings of
\cite[See ][, p.238]{balardy2022empirical}. During the IDAs the
spread increases significantly, which is expected given that the
spread is negatively correlated to volume according to XX. The median
spread of the average trading session is 0.89 €/MWh, which is
considerably less than what was found in previous studies. For
comparison, the median spread in 2015 reported by Balardy was about
3.52 €/MWh. This decrease can be attributed to the annually
increasing volume in the intraday market. Although the median spread
is more than three times smaller than the median spread in 2015, it
is still 89 times the tick size, which compared to equity markets is
massive. In equity markets the spread is typically few times the tick
size at most. The weekend effect in liquidity can also be observed
through the spread as Balardy reports in that the spread is on
average 13\% higher during the weekends. The tick size also has a
large effect on the spread distribution. Given that the tick size of
0.01 €/MWh is quite small compared to the average price per megawatt,
the observed distribution is not that surprising.
% As idea here, make a small table comparing other studies bid-ask spreads
Not surpising, the spread is also connected to volumes, as the more
frequent the trades, the better the spread, and also to the
volatility of the spread itself. More frequently traded assets tend
to have less volatile quoted spreads \cite{balardy2022empirical}. A
slightly older study conducted by \cite{hagemann2013empirical} found
a lower mean bid-ask spread of about 2.97€/MWh. As the durations of
the trading sessions per product differ in their lengths, a smooting
of the bid-ask spread can be observed, the longer the trading session
lasts \cite[See ][, p.238]{balardy2022empirical}. Concluding, the
bid-ask spread generally follows more like an L-shape instead of a
U-shape pattern as is usually observed in other financial markets
\cite[See ][, p.238]{balardy2022empirical}. As for the main drivers
of the bid-ask spread, increased volatility positively affects the
bid-ask spread, especially pronounced for off-peak produtcs
\cite[See][, p.246]{balardy2022empirical}. Not surprisingly,
fundamental variables like load, wind, solar and their associated
forecast errors (positive and negative) have a negative impact for
the bid-ask spread \cite[See][, p.247]{balardy2022empirical}. The
Bid-ask spread has two important dependency relationships: First, the
impact of the Bid-ask spread on posted LO depths can't be ignored,
since the bid-ask spread is rarely tick size and often exceeds it by
several magnitudes.\footnote{\cite[See ][, p.78]{luckner2021hawkes}}
Second, it is relevant for the arrivals of MO, as we have seen that
they appear in situations where the spread is favorable for them and
narrow.
%------------------------------------------------------------------------------------
%	                  FIGURE - START
%------------------------------------------------------------------------------------
\begin{figure}[htbp]
    \vspace*{0.5cm}
    \centering
    \blendfigure{1.0\textwidth}{9cm}{intraday_spread_pattern_iqr.png}
    \captionfootcite{Intraday spread pattern (IQR).}{Based on \cite[See ][, p.75]{cartea2015algorithmic}.}\label{fig:intraday_spread_pattern_iqr}
\end{figure}
%------------------------------------------------------------------------------------
%	                  FIGURE - END
%------------------------------------------------------------------------------------
Pictured in the figure above, quoted spreads are initially high,
decline rapidly during the first half-hour of trading, and are mostly
constant throughout the remainder of the day until the last (half)
hour of trading, when the spread rapidly declines. This pattern in
quoted spreads is also seen in the effective
spread.\footnote{\cite[See ][, p.75]{cartea2015algorithmic}} We can
resonate that during a period with a constant flow of information
more trading and lower spreads occur together.\footnote{\cite[See ][,
        p.75]{cartea2015algorithmic}}
% Result for modelling: Bid-ask spread is also dependent on the time-to-delivery, so add that as state variable
% Result for modelling: Impact of Bid-ask spread on posted LO depths can't be ignored, since the bid-ask spread is rarely tick size and often exceeds it by several magnitudes. So add that as state variable
% Result for modelling: Bid-ask spread also importnat for MO arrivals
% Result for modelling: Bid-ask spread is negatively correlated with volume and positively correlated with volatility, so add that as state variable
%------------------------------------------------------------------------------------
\subsection{Market Depths and Trade Size}\label{sub_sec:market_depths_and_trade_size}
%------------------------------------------------------------------------------------
Confirmed here as already presented in
\hyperref[sec:market_making_problem]{2.3 The Market Making Problem},
market depths and associated trade sizes are not independent of
another.\footnote{\cite[See ][, p.80]{cartea2015algorithmic}} Depth
in this context refers to the available volumes resting on the price
levels in the LOB at a given time, ready for immediate
execution.\footnote{\cite[See ][, p.79]{cartea2015algorithmic}} This
behaviour can be motivated as market participants look to minimize
their market impact, splitting larger block trades into several
smaller ones. This indicates two phenomens: First, as already stated,
the depth of the LOB and size of the incoming orders are highly
correlated. But also secondly, this splitting induces, and at least
in parts, explains the already observed clustering of MO arrivals. To
get an intuition for the resting depths of the LOB, the figure below
shows that the available liquidity on both sides is relatively
similar and centered around 5 megawatt quantities. Five megawatts is
also the minimum visible quantity for an iceberg order on EPEX SPOT,
which might explain the high number of orders with such a quantity.
Interestingly, the volume on both sides drops visibly after the third
IDA. This is most likely caused by execution algorithms getting more
aggressive towards the end of the trading session as they split
orders into small pieces and try to execute passively in order to get
a better price. Verifying this is of course not possible since all
order information is anonymous.
%------------------------------------------------------------------------------------
%	                  FIGURE - START
%------------------------------------------------------------------------------------
\begin{figure}[htbp]
    \vspace*{0.5cm}
    \centering
    \blendfigure{1.0\textwidth}{9cm}{lob_depth_top_of_book.png}
    \captionfootcite{Top of the book depth as a function of time to delivery.}{Based on \cite[See ][, p.75]{cartea2015algorithmic}.}\label{fig:lob_depth_top_of_book}
\end{figure}
%------------------------------------------------------------------------------------
%	                  FIGURE - END
%------------------------------------------------------------------------------------
Looking beyond just the best bid and ask levels, the figure below
shows the depths for the best five price levels on both sides. In
combination with the previous figure, it is evident, that the number
of individual orders at each price level is almost always one or two,
which means that the order queues are very weak. For market making
purposes strong queues would be optimal since then the price impact
and the resulting negative effect on the market makers inventory
would be smaller. Figure 3.5, which shows the cumulative liquidity
and the total number of the orders at the top five levels of both
sides of the LOB, confirms the earlier observation that towards the
end of the trading session the orders get smaller and the liquidity
at the top levels of the book gets more scarce.
%------------------------------------------------------------------------------------
%	                  FIGURE - START
%------------------------------------------------------------------------------------
\begin{figure}[htbp]
    \vspace*{0.5cm}
    \centering
    \blendfigure{1.0\textwidth}{9cm}{lob_depth_top5_levels.png}
    \captionfootcite{LOB depth for top 5 price levels as a function of time to delivery.}{Based on \cite[See ][, p.75]{cartea2015algorithmic}.}\label{fig:lob_depth_top5_levels}
\end{figure}
%------------------------------------------------------------------------------------
%	                  FIGURE - END
%------------------------------------------------------------------------------------
In the figure below, we plot the mean and median distance between the
top price levels over the average trading session. The distribution
for both sides is very similar to the bid-ask distribution, excluding
the IDAs. The distances get smaller towards delivery which would
indicate that the total liquidity in the LOB increases. During the
IDAs and the final hour of trading, the distances jump up as a
natural result of the XBID being turned off.
%------------------------------------------------------------------------------------
%	                  FIGURE - START
%------------------------------------------------------------------------------------
\begin{figure}[htbp]
    \vspace*{0.5cm}
    \centering
    \blendfigure{1.0\textwidth}{9cm}{mean_median_dist_5_price_levels.png}
    \captionfootcite{Mean and median distance between the top 5 price levels as a function of time to delivery.}{Based on \cite[See ][, p.75]{cartea2015algorithmic}.}\label{fig:mean_median_dist_5_price_levels}
\end{figure}
%------------------------------------------------------------------------------------
%	                  FIGURE - END
%------------------------------------------------------------------------------------
As the arrivals of MO depend on the Bid-ask spread in favor for
tighter conditions, we can also conclude that the posted LO depths
are dependent as well. Market depths correlate negatively with higher
bid-ask spreads \cite[see ][, p.237]{balardy2022empirical}, what is
especially pronounced for off-peak hours and weekeends. Higher
bid-ask spreads primarily result from higher sell prices and lower
buy prices \cite[][, p.237]{balardy2022empirical}, leaving business
days and peak-contracts as the most liquid combinations in the
market. As the spreads widen, the MM tend to post their LOs closer to
the midprice, which is also a natural reaction to the lower liquidity
in the LOB.
% Result for modelling: market depth and order sizes are not independent
% Result for modelling: market depth very low with only one or two resting LO per price level -> high market impact due to order walking the book
% Result for modelling: market depth and correlate negatively with higher bid-ask spreads
% Result for modelling: Not sufficient to only model best bid and ask but rather also the 1-3 other price levels
%------------------------------------------------------------------------------------
\subsection{Market Impact}\label{sub_sec:market_impact}
%------------------------------------------------------------------------------------
As we have seen from the previous market-depth analysis, for the CID
market, the market impact plays a big role since resting LO volumes
at the best bid and ask are rather small and would be entirely
consumed quickly by large incoming MO, depleting these price levels,
`walking the book' and attaining a stepwise worsening average price.
Price impact measures the cost of executing larger trades via their
impact of trading on prices.\footnote{\cite[See
    ][,p.88]{cartea2015algorithmic}} A linear relationship between price
changes and net order flow is plausible for a wide range of midprice
changes, what justifies the use of a simple linear model to measure
the market impact. By the simple linear model, we can see that there
is a positive relationship between net order flow and price changes.
That is when net order flow is positive (negative), that is more
(less) buy than sell MO, the midprice tends to increase
(decrease).\footnote{\cite[See ][,p.88]{cartea2015algorithmic}} The
price impact can be further decomposed into two components: the
permanent and the temporary impact. An exemplary figure of the
results for the temporary price impact can be seen below:
%------------------------------------------------------------------------------------
%	                  FIGURE - START
%------------------------------------------------------------------------------------
\begin{figure}[htbp]
    \vspace*{0.5cm}
    \centering
    \blendfigure{1.0\textwidth}{9cm}{order_flow_mid_price_drift.png}
    \captionfootcite{Order Flow and effect on the drift of midprice of INTC.}{Based on \cite[See ][, p.75]{cartea2015algorithmic}.}\label{fig:order_flow_mid_price_drift}
\end{figure}
%------------------------------------------------------------------------------------
%	                  FIGURE - END
%------------------------------------------------------------------------------------
The third picture in the Figure shows how the slope of this linear
impact model fluctuates throughout the entire day. We see that the
largest impact tends to occur in the morning, then this impact
flattens and stays flat throughout the day, and towards the end of
the day it lessens. Such a pattern is seen in a number of assets and
is consistent with the reduction in spreads and increases in depth we
have documented earlier. So, our very preliminary and limited
analysis suggests that order arrival seems to be followed by further
order arrival on both sides of the book, and more on its own side
than on the other side of the book. As we will now see, MOs do not
arrive at random times. They tend to arrive when spreads are narrow,
and opportunistically hit orders that are posted closer to the mid
price, so it is only natural that we should observe a wider spread
after an execution.\footnote{\cite[See
    ][,p.87]{cartea2015algorithmic}}
%------------------------------------------------------------------------------------
%	                  FIGURE - START
%------------------------------------------------------------------------------------
\begin{figure}[htbp]
    \vspace*{0.5cm}
    \centering
    \blendfigure{1.0\textwidth}{9cm}{market_impact.png}
    \captionfootcite{Market impact.}{Based on \cite[See ][, p.88]{cartea2015algorithmic}.}\label{fig:market_impact}
\end{figure}
%------------------------------------------------------------------------------------
%	                  FIGURE - END
%------------------------------------------------------------------------------------
% Result for modelling: Market impact overall relevant as market depths are very thin
% Result for modelling: Time dependent market impact
% Result for modelling: Linear Market impact is justified and shall be used
% Result for modelling: Positive Relationship between net order flow and price changes, so include it as a state variable
%------------------------------------------------------------------------------------
\subsection{Messages and Cancellations}\label{sub_sec:messages_and_cancellations}
%------------------------------------------------------------------------------------
Key components forming the shape of the LOB are the arrival and
cancellation rates of LOs and MOs. As these are critical inputs for
the overall optimisation problem of the MM, in the sense that their
intensities determine the frequency of the LOs being lifted and the
resulting gain from the spread. Descriptive statistics are provided
in the table below:
%------------------------------------------------------------------------------------
%	                  FIGURE - START
%------------------------------------------------------------------------------------
\begin{figure}[htbp]
    \vspace*{0.5cm}
    \centering
    \blendfigure{1.0\textwidth}{9cm}{num_messages_per_day.png}
    \captionfootcite{Number of messages per day.}{Based on \cite[See ][, p.91]{cartea2015algorithmic}.}\label{fig:num_messages_per_day}
\end{figure}
%------------------------------------------------------------------------------------
%	                  FIGURE - END
%------------------------------------------------------------------------------------
Analysing the arrival rates of MOs and LOs, Figure 3.7 shows the
average intensity of MO arrivals during a trading session for hour 14
delivery.
%------------------------------------------------------------------------------------
%	                  FIGURE - START
%------------------------------------------------------------------------------------
\begin{figure}[htbp]
    \vspace*{0.5cm}
    \centering
    \blendfigure{1.0\textwidth}{9cm}{mo_arrival_rates.png}
    \captionfootcite{MO arrival distribution as a function of time to delivery.}{Based on \cite[See ][, p.91]{cartea2015algorithmic}.}\label{fig:mo_arrival_rates}
\end{figure}
%------------------------------------------------------------------------------------
%	                  FIGURE - END
%------------------------------------------------------------------------------------
The distinction between aggressive and non-aggressive MOs is that
aggressive MOs move the current mid-price i.e. the volume of the MO
is large enough to deplete all the liquidity at the current best bid
or ask, and non-aggressive MOs have less volume than is currently
available at the best bid or ask. The intensity of MO arrivals is
highly correlated with the volume pattern in Figure 3.1, but the
pattern is smoother and follows an exponential shape. The intensity
drops significantly only in the final hour of trading, and the IDAs
seem to have little to no impact on the intensity. The share of
aggressive MOs to non-aggressive MOs seems to be constant. Von
luckner and Kiesel also provide strong empirical evidence about the
clustering of MOs in the German intraday market in. Clustering means
that successive orders are more likely to be of the same sign, which
can be attributed to market participants splitting their MOs into
smaller pieces, following the behavior of other traders or acting on
the same information in timely proximity. The ability to cancel and
place new limit orders is critical for the market maker in order to
change their LO placement in the LOB upon arrival of new market
events. Figures 3.8 and 3.9 present the intensities of LO arrivals
and cancellations for the hour 14 delivery. Aggressive LO arrival is
defined as a LO that arrives inside the spread thus changing the
mid-price and non-aggressive arrival is an order that arrives at the
current best bid or ask, or a worse price. Aggressive cancellation is
defined as a canceled order that moves the mid-price which happens
when the canceled order is the only order at the best bid or ask. The
patterns for LO additions and cancellations are almost identical,
which is not a coincidence: most of the orders are placed and then
canceled in the next few seconds. The patterns for additions and
cancels are not as smooth as the MO arrival pattern and are
significantly affected by the IDAs. During the IDAs, the intensities
for both additions and cancellations drop significantly and jump
right after when XBID is turned on again. The drop in intensity in
the final hour is also very clear. The share of aggressive LO
additions and cancels is relatively constant with small spikes during
the IDAs due to lower liquidity in the LOB.
%------------------------------------------------------------------------------------
%	                  FIGURE - START
%------------------------------------------------------------------------------------
\begin{figure}[htbp]
    \vspace*{0.5cm}
    \centering
    \blendfigure{1.0\textwidth}{9cm}{lo_ca_arrival_rates.png}
    \captionfootcite{LO and CA arrival distribution as a function of time to delivery.}{Based on \cite[See ][, p.91]{cartea2015algorithmic}.}\label{fig:lo_ca_arrival_rates}
\end{figure}
%------------------------------------------------------------------------------------
%	                  FIGURE - END
%------------------------------------------------------------------------------------

%------------------------------------------------------------------------------------
%	                  FIGURE - START
%------------------------------------------------------------------------------------
\begin{figure}[htbp]
    \vspace*{0.5cm}
    \centering
    \blendfigure{1.0\textwidth}{9cm}{messages_by_dist_to_mid.png}
    \captionfootcite{Messages by distance to midprice.}{Based on \cite[See ][, p.93]{cartea2015algorithmic}.}\label{fig:messages_by_dist_to_mid}
\end{figure}
%------------------------------------------------------------------------------------
%	                  FIGURE - END
%------------------------------------------------------------------------------------
From the above table we find that the distance to the midprice is
(naturally) shorter, and we can compute that 36\% of orders were
executed at prices between 1 and 3 cents from the midprice, and 42\%
at prices between 4 and 6, so that only 22\% of executions were
relatively far from the midprice (7+ cents away).\footnote{\cite[See
    ][, p.93]{cartea2015algorithmic}}
%------------------------------------------------------------------------------------
%	                  FIGURE - START
%------------------------------------------------------------------------------------
\begin{figure}[htbp]
    \vspace*{0.5cm}
    \centering
    \blendfigure{1.0\textwidth}{9cm}{surv_func_dist_from_mid.png}
    \captionfootcite{Survival function of distance from midprice.}{Based on \cite[See ][, p.94]{cartea2015algorithmic}.}\label{fig:surv_func_dist_from_mid}
\end{figure}
%------------------------------------------------------------------------------------
%	                  FIGURE - END
%------------------------------------------------------------------------------------
In the above figure of total executions, as the distance from the
price at which the original LO was posted increases. This represents
an approximation to the 'fill probability' - the probability that a
posted order is executed, what we will also later use in the
modelling. Assuming a certain market share of $\rho$ that the MM of
the total traded volume has, or wishes to achieve, and that is chosen
in advance, we can use this directly in the sense
of\footnote{\cite[See ][, p.9]{law2015pure}} multiplying the
effective arrival rate of market orders hitting his limit orders by
$\rho$. For example, when $\rho$=5\%, that means the effective
arrival rate of market orders hitting his limit orders is around 5\%
of all market orders. He may achieve this by continuously adjusting
his limit order quantity to be 5\% of the total quantity in the queue
but the detailed mechanism is outside the scope of this
study.\footnote{\cite[See ][, p.61]{law2015pure}} The market share
$\rho$ will be multiplied to the cash and quantity dynamics of the
market maker while the price will evolve according to all orders in
the market. When he has limit orders in the bid queue, the increase
in quantity will be the number of shares v of the market order
hitting the bid multiplied by the market share $\rho$ (assume market
maker’s limit orders are distributed evenly within the
queue).\footnote{\cite[See ][, p.62]{law2015pure}}

This indicates that market buy orders tended to occur much closer to
the midprice than market sell orders on this particular day, which
had an overall positive order flow for AAPL shares and a slight price
increase from market open to market close. In Figure 4.14, we have
also included total executions separated by the time of day: the
first half hour after the market opens (Mkt Start), the last half
hour before the market closes (Mkt Close), and the time in between
(Intraday). We observe that Mkt Close tends to be below that of
Intraday, implying that during the last half hour of trading,
executions tend to be close to the midprice, which is consistent with
the pattern of the quoted spread in Figure 4.6.
% Result for modelling: most of executed postings very close to midprice -> Only try to model levels +-3 from midprice or best bid and best ask
% Result for modelling: definitely filter out the times of the IDA as this will negatively impact and bias the intensity estimations
% Result for modelling: Clustering of MO arrivals
% Result for modelling: Consider the best three levels for modelling as there is the majority of activity and volumes
%------------------------------------------------------------------------------------
\section{Empirical Intraday activity patterns and stylized facts}\label{sec:empirical_patterns_and_stylized_facts}
%--------------------------------------------------------------------------------
As we have seen in the preceeding empirical analysis, the data
exhibits the following key charactersistics, that set the stage for
our benchmarking framework for the mathematical modelling and its
performance assessment:\footnote{\cite[See ][, p.13]{law2015pure}}
% Maybe make a table out of that where I list all the papers of the power trading sector
% that confirm the single finding in their research
\begin{enumerate}
    \item irregular time intervals between arrivals
    \item discrete state space of price ticks and volume lot sizes
    \item intraday seasonality (more activities around market open and close)
    \item arrival clustering
    \item self-excitation from its own history
    \item cross-excitation from other order types
    \item cross-excitation from the history of other assets
    \item long memory of excitation effect
    \item sparse order book characteristsics (low liquidity, low market depth,
          wide bid-ask spread)\footnote{\cite[See ][,
                  p.236]{balardy2022empirical}}
    \item Non-Markovian dynamics (the future evolution of the LOB depends not
          only on the current state but also on the past history of order
          arrivals and executions)\footnote{\cite[See ][,
                  p.52]{cartea2015algorithmic}}
\end{enumerate}

As with the newest regulatory requirements, the todays standard will
be the 15-minute contracts, especially item seven proposes a
complexity increasing hurdle. Suspectively the closer a neighbouring
contract is, the more cross-excitation impact will be observed,
suggesting few contracts before and after the current contract, as
the most relevant ones for the modelling. We chose to ignore this
complexity for the sake of simplicity and focus on the single
contract modelling.
%------------------------------------------------------------------------------------
\chapter{Mathematical Modelling}\label{ch:mathematical_modelling}
%------------------------------------------------------------------------------------
% Centers around Order Flow Modelling
In order to realistically model the LOB dynamics and derive the
optimal market making strategy, we need to choose a suitable
mathematical framework that can replicate the key empirical
characteristics of the LOB. A natural framework is offered by the
Hawkes-Process family, which is a class of self-exciting point
processes that can capture the clustering and dependence structure of
order arrivals in the LOB.\footnote{Note: In theory the order arrival
intensities can be any predictable non-negative stochastic processes, \cite[see ][, p.39]{korhonen2025market}.} By also including cross-excitation among
different order types, this process family also offers modelling
capabilities for complex interactions.

Explicitly not assuming that the MM can adjust bid and ask prices
continuously, the MM's decision problem is formulated as a form of an
optimal impulse-control switching problem, see \cite{law2015pure}.

Not only is this assumption not realistic but also not properly
penalised in the existing discussions, as altering order parameters
means losing the queue position due to the applied price-time
priority principle most exchanges operate on.

In the following sections, we will first introduce the basic concepts
of Hawkes processes, as member of the general point process family,
and then describe how we can use this framework to model the order
arrivals and mid-price dynamics in the electricity markets,
ultimately leading to the derivation of the optimal market making
strategy. A complete taxonomy of the main process families belongig
to the point processes is given below.
%------------------------------------------------------------------------------------
%	                  FIGURE - START
%------------------------------------------------------------------------------------
\begin{figure}[htbp]
    \vspace*{0.5cm}
    \centering
    \blendfigure{1.0\textwidth}{7cm}{point_process_family.png}
    \captionfootcite{Shcematic overview of the point process family.}{Based on \cite{kiesel_efc19_2019}, own illustration.}\label{fig:point_process_family}
\end{figure}
%------------------------------------------------------------------------------------
%	                  FIGURE - END
%------------------------------------------------------------------------------------

%------------------------------------------------------------------------------------
\section{The Point Process Family}\label{sec:point_process_family}
%------------------------------------------------------------------------------------
Starting with the basic concepts of the general point process family,
we can define a point process as a random collection of points in
time, where each point represents the occurrence of an event. This
basic model for representing event streams assumes that their
occurrence is random, i.e. independent and can be described by a
probability
distribution.\footnote{\cite{palm1943intensitatsschwankungen}}
%------------------------------------------------------------------------------------
\subsection{Homogeneous Point Processes}\label{sub_sec:homogeneous_point_processes}
%------------------------------------------------------------------------------------
\cite[See ][, p.3]{oxford_hawkes_2020}
\begin{definition}[Point process]\label{def:point_process}
    A point process is a stochastic process $(T_k)_{k\in\mathbb{N}}$
    taking values in $\mathbb{R}_{+}$ such that for every $k\in\mathbb{N}$, $T_k \leq T_{k+1}$.
\end{definition}

\begin{definition}[Counting Process]\label{def:counting_process}
    Given a point process $(T_k)_{k\in\mathbb{N}}$, its counting process is defined by
    \begin{equation}\label{eq:counting_function}
        N(t)=\sum_{k\in\mathbb{N}}\mathbf{1}_{\{T_k\le t\}}
    \end{equation}
    and counts the number of events that have occurred up to time $t$. The
    process is called \emph{simple} if no two events occur at the same time,
    i.e., $\mathbb{P}(T_k=T_j)=0$ for all $k\neq j$.
\end{definition}
In reality of course, multiple incoming messages can have the exact same assigned timestamp,
suggesting a \emph{non-simple} point process. However, for the sake of simplicity,
we will assume a simple point process in our modelling framework, as these cases are very
rare and only make up x \% of observations (timestamped on nanoseconds).\footnote{Assume no two types
    of orders can arrive at the exact same time, if the exchange operates multiple
    servers it can of course accept two incoming order types at the same instant, \cite[see ][p. 45]{law2015pure}}

An important characteristic of the counting process is its so-called
\emph{conditional intensity}:

\begin{definition}[Conditional Intensity]\label{def:conditional_intensity}
    The conditional intensity $\lambda(t)$ of a counting process $N(t)$ is defined as the limit
    \begin{equation}\label{eq:conditional_intensity}
        \lambda(t) := \lim_{\Delta t \to 0} \frac{\mathbb{P}(N(t+\Delta t)-N(t)=1|\mathcal{F}_t)}{\Delta t}
    \end{equation}
    where $\mathcal{F}_t$ is the filtration generated by the counting process up to time $t$,
    i.e. $\mathcal{F}_t : = \sigma(N(s): 0 \leq s \leq t)$.
\end{definition}

Another important quantity in the study of point processes is the
point-wise integral of the conditional intensity of the process,
$\Lambda(t)$, i.e.
\begin{definition}[Compensator]\label{def:compensator}
    \begin{equation}\label{eq:compensator}
        \Lambda(t) := \int_0^t \lambda(s) ds
    \end{equation}
    We call this value the compensator of the process, and we have the
    following useful result about compensators (Lemma X.Y.Z in XX).
\end{definition}
% This is the regular compensator for martingales (Doob Decomposition)
\begin{theorem}[$\mathcal{F}$-Martingale]\label{thm:f_martingale}
    Suppose the $\mathcal{F}_t$-adapted counting process $(N(t))_{t \in \left[0, \infty \right)}$ has a left-continuous
    intensity $\lambda(t)$. Then the process $M(t) := N(t) - \Lambda(t)$
    is an $\mathcal{F}_t$-martingale, i.e. for every $0 \leq s \leq t$,
    \begin{equation}\label{eq:f_martingale}
        \mathbb{E}[M(t) | \mathcal{F}_s] = M(s).
    \end{equation}
    Counting processes preserve three notions of regularities, namely:
    \emph{stationarity}, \emph{simplicity} and \emph{orderliness}.\\
    A counting process $N$ is said to be (strictly) \emph{stationary} if $N(t + s) - N(s)$
    has the same distribution for all $s \geq 0$. $N$ is said to be \emph{simple} when:
    \begin{equation}
        \mathbb{P}(N({t}) = \textit{0 or 1 for all t}) = 1
    \end{equation}
    and orderly, wehn it is first stationary and:
    \begin{equation}
        \mathbb{P}(N(t+h) - N(t) \geq 2) = o(h) \textit{as h} \rightarrow 0
    \end{equation}.
\end{theorem}
One of the most prominent members of counting processes is the \emph{Possion Process},
that can be defined in serveral ways.
Based on its intensity it may be defined as:
\begin{definition}[Possion Process]\label{def:possion_process}
    Let $N$ be a counting process and consider a sequence $\left(s_{i},{} t_{i}\right]$
    of half-open intervals with $0 \leq s_{i} < t_{i} \leq s_{i+1}$, $i \in \{1, \ldots, k\}$.
    $N$ is called a Poisson process with intensity $\lambda > 0$ if:
    \begin{equation}
        \mathbb{P}[N(t_{i}) - N(s_{i}) = n_{i}, i \in {1, \ldots, k}] = \prod_{i=1}^{k} \frac{(\lambda(t_{i}-s_{i}))^{n_{i}}}{n_{i}!} e^{-\lambda(t_{i}-s_{i})}
    \end{equation}

    It follows from this definition that a Possion Process satisfies the
    following conditions:
    \begin{enumerate}
        \item The number of events in an interval $\left[s, t\right),~0 \leq s <
                  t$, follows a Poisson distribution with mean $\lambda (t-s)$
        \item the process has \emph{independent} increments, i.e. the number of
              arrivals in disjoint intervals are independent.
    \end{enumerate}
\end{definition}
Allowing the intensity $\lambda$ to vary with time, we transition from
the homogeneous to the non-homogeneous Possion Process.\footnote{\cite[See ][, p.4]{oxford_hawkes_2020}}
%------------------------------------------------------------------------------------
\subsection{Non Homogeneous Point Processes}\label{sub_sec:non_homogeneous_point_processes}
%--------------------------------------------------------------------------------
In a non-homogenous Poisson process, the (infinitesimal) probability
of an event happening at time t may vary with t, but it is still
independent of other events.\footnote{\cite[See ][,
        p.1]{mei2017neural}}

An importnat characteristic of the Possion Process is its
\emph{memorylessness}, i.e. the distribution of future arrival times
only depends on information available at the current time $t$, and
not what happened in the past leading to $t$.\footnote{\cite[See ][,
        p.4]{oxford_hawkes_2020}} To adjust for this \emph{memoryless}
property, a historic dependency structure has to be introduced,
leading to the \emph{Hawkes-Process}.
%------------------------------------------------------------------------------------
\subsection{One Dimensional Hawkes Processes}\label{sub_sec:one_dimensional_hawkes_processes}
%------------------------------------------------------------------------------------
A Hawkes process\footnote{\cite[See ][, p.]{hawkes1971spectra}}
(Liniger, 2009) supposes that past events can temporarily raise the
probability of future events, assuming that such excitation is
positive, additive over the past events, and exponentially decaying
with time, see \hyperref{thm:f_martingale}{Theorem
    1}.\footnote{\cite[See ][, p.1]{mei2017neural}} Real world scenarios
often violate these assumptions, as an event can also be inhibitious
rather than excitatory, and the excitation effect may not be strictly
additive over time as effect strengths might dampen down or that an
exponentially decay can't be guaranteed.\footnote{\cite[See][,
        p.2]{mei2017neural}} The output of a Hawkes process is essentially an
event intensity, which is an instantantaneous
probability.\footnote{\cite[See][, p.2]{mei2017neural}}

\begin{definition}[Hawkes Process]\label{def:hawkes_process}
    A Hawkes process is a simple point process with a conditional intensity function $\lambda(t)$
    defined as:
    \begin{equation}\label{eq:hawkes_intensity}
        \lambda(t) = \mu + \int_{t_i < t} \phi(t - s) dN_{s}
    \end{equation}
    where $\mu \in \mathbb{R}^+$ is the baseline intensity of the process,
    $\phi(t): \mathbb{R}^+ \to \mathbb{R}^+$ is the excitation kernel function, and $t_i$ are the times of past events.
\end{definition}\footnote{\cite[See ][, p.5]{oxford_hawkes_2020}}

In dimension $d = 1$, the process exhibits \emph{self-excitation},
where past event occurences, spark an increase in the likelihood of
observing events in the near future. Note that in $d=1$, only one
specific type of event is considered for the self-excitation. It is
of course possible and rather plausible, that in realistsic
scenarios, multiple types of events are responsible for future
excitations, proposing the need to incporporate a dependency
structure for mutual event type impacts.\footnote{\cite[See ][,
        p.5]{oxford_hawkes_2020}}
% This is the exact reason why we need a d-dimensional framework, since for a specific considered event
% type, like MO arrivals, excitation of additional future MO arrivals also depends on the other queues like LO arrivals, Cancellations
To motivate a later extension, in effect, Hawkes processes and
continuous-time Markov chains can have either an event viewpoint or a
state viewpoint, respectively.\footnote{See \cite{morariu2018hybrid}
    p.2}
%------------------------------------------------------------------------------------
\subsection{Multi Dimensional Hawkes Processes}\label{sub_sec:multi_dimensional_hawkes_processes}
%------------------------------------------------------------------------------------
\cite[See ][,p.5]{oxford_hawkes_2020}:

\begin{definition}[Multivariate Hawkes Process]\label{def:multivariate_hawkes_process}
    A $d$-dimensional counting process $N = (N_1, \ldots, N_d)$, is a
    multivariate Hawkes process when:
    \begin{equation}\label{eq:multivariate_hawkes_intensity}
        \lambda_{i}(t) = \mu_{i} + \sum_{j = 1}^{d} \int_{0}^{t} \phi_{ij}(t - s) dN_{j, s}
    \end{equation}
    where $\mu_{i} \in \mathbb{R}^d$ are the baseline intensities, $t_i$ are the times of past events and
    $\phi_{ij}(t): \mathbb{R}^+ \to \mathbb{R}^{d \times d}$ are the excitation kernel functions,
    and satisfy:
    \begin{enumerate}
        \item $\phi_{ij}(t) \geq 0$ for all $t \in \mathbb{R}$ and $i, j \in \{1, \ldots, d\}$
        \item $\phi_{ij}(t) = 0$ for all $t \leq 0$ and $i, j \in \{1, \ldots, d\}$
        \item $\phi_{ij}(t)$ is locally integrable for all $i, j \in \{1, \ldots, d\}$ i.e. $\phi_{ij}(t) \in L_{1, loc}(\Omega)$
    \end{enumerate}
\end{definition}

As in XX, two alternative notations of
\hyperref[eq:multivariate_hawkes_intensity]{Equation 4.9} can be
considered, allowing for a more compact notation and extended
interpretation of its components. The first one, being the
\emph{convolutional} notation:
\begin{equation}
    \Lambda(t) = \mu + \Phi * dN_t
\end{equation}
where $\Phi * dN_t$ is the convolution of the kernel function $\Phi$ with the counting process $dN_t$.
The second one, being the \emph{matrix} notation:
\begin{equation}
    \lambda(t) = \mu + \int_{0}^{t} \Phi(t - s) dN_{s}
\end{equation}
%------------------------------------------------------------------------------------
\subsection{Adding State Dependencies and Marks}\label{sub_sec:adding_state_dependencies_and_marks}
%------------------------------------------------------------------------------------
After having discussed the one- and d-dimensional Hawkes Process, the
state dependent and marked variants should be introduced.
Understanding how these two properties are essential for a realistic
modelling of the LOB dynamics, we must first differentiate them along
two dimensions:
\begin{enumerate}
    \item where information originates from and
    \item how it affects the intensity ($\lambda_{t}$) of future order
          arrivals.
\end{enumerate}

Adding state dependency and marks is also motivated by the findings
here: \cite[see ][, p.83]{cartea2015algorithmic} where we clearyl see
that the size of an order (so its magnitude) impacts the arrival
rates of after following orders heavily. Here it is alos suggested
that order arrivals seem to be followed by further order arrivals on
both sides of the book, and more on its own side than on the opposing
side \cite[see ][, p.83]{cartea2015algorithmic}.

%------------------------------------------------------------------------------------
\subsubsection{The Marked Hawkes Process (MHP)}\label{sub_sub_sec:the_marked_hawkes_process}
%--------------------------------------------------------------------------------
The Marked Hawkes Process can be thought of as containing
Event-Driven Info. The intensity changes based on a specific property
attached directly to an incoming event at the exact millisecond it
arrives. An incoming order has attributes like its size, direction,
and price level, which are all relevant for the future order arrival
intensity. The Marked Hawkes Process is therefore a natural extension
of the standard Hawkes Process, where the intensity is only a
function of the past event times.

\begin{paragraph}{Mathematical formulation}
    Every time an event \(i\) occurs at time
    \(t_{i}\), it carries a "mark" (\(m_{i}\)). In market microstructure,
    this mark is almost always the Order Size / Volume or Trade
    Direction:
    \begin{equation}
        \lambda_{t} = \mu + \sum_{t_{i}<t} \phi(t-t_{i}, m_{i})
    \end{equation}
\end{paragraph}

The intuition behind including marks is that the impact of an event
on future order arrivals is not only determined by its occurrence but
also by its characteristics. For example, a large market order
(carrying the mark $m_{i}$) may have a more significant impact on the
afterfollowing excitation ripple (a bigger shock through the kernel
function: \(\phi(t-t_{i}, m_{i})\)) and thus a higher intensity of
future orders than a small MO, especially considering situations
where the Market-depth is thin. In this thesis, we will explore if
incorporating the \emph{Volume} of an order as mark will carry any
beneficial information for the future order arrival intensity.

Memorywise: The system only looks backward at the history of
individual events and their specific sizes.
%------------------------------------------------------------------------------------
\subsubsection{The State Dependent Hawkes Process (SDHP)}\label{sub_sub_sec:the_state_dependent_hawkes_process}
%--------------------------------------------------------------------------------
The State-Dependent Hawkes Process represents Environment-Driven
Info. The intensity changes based on the ambient state of the entire
system at time \(t\), regardless of which specific event caused that
state.

\begin{paragraph}{Mathematical formulation}
    Introduced to fix a major flaw in pure
    Hawkes models: traditional Hawkes processes ignore the geometry of
    the order book. In a state-dependent model, the kernel itself is
    modulated by a piecewise constant state variable, \(X_{t}\),
    representing the continuous state of the LOB.

    \begin{equation}
        \lambda_{t} = \mu(X_{t}) + \sum_{t_{i}<t} \phi_{X_{t_{i}}}(t-t_{i}, X_{t})
    \end{equation}
\end{paragraph}

The intuition behind including state dependencies is that the state
\(X_{t}\) could be the bid-ask spread or the queue imbalance. If the
spread is at its minimum tick size, the probability of a
price-improving limit order dropping in is physically zero. The model
dampens or accelerates intensities based on the geometry of the book
layout when the order hits.

Memorywise: The system is fully coupled. The arrival of an order
changes the book state (\(X_{t}\)), and the new book state instantly
updates the baseline (\(\mu \)) and excitation physics (\(\Phi \)) of
future orders.

%------------------------------------------------------------------------------------
\subsubsection{The Unified Approach: The State dependent marked Hawkes Process (SDMHP)}\label{sub_sub_sec:the_unified_approach}
%--------------------------------------------------------------------------------
In a combined model, the intensity \(\lambda _{t}\) of order arrivals
depends simultaneously on the changing geometry of the book
(\(X_{t}\)) and the specific properties of past historical trades
(\(m_{i}\)).

\begin{paragraph}{Mathematical formulation}
    \begin{equation}
        \lambda_{t} = \mu(X_{t}) + \sum_{t_{i}<t} \phi_{X_{t_{i}}}(t-t_{i}, m_{i}, X_{t})
    \end{equation}
    with the following components:
    \begin{itemize}
        \item \(X_{t}\): The state of the LOB at time \(t\), which can include variables such as the bid-ask spread, queue imbalance, and other relevant market conditions.
        \item \(m_{i}\): The mark associated with the \(i\)-th event, which can represent characteristics such as order size or trade direction.
        \item \(\mu(X_{t})\): The baseline intensity, which is a function of the current state of the LOB, capturing the underlying rate of order arrivals based on the current market conditions.
        \item \(\phi_{X_{t_{i}}}(t-t_{i}, m_{i}, X_{t})\): The coupled excitation kernel, which depends on the time since the last event \(t-t_{i}\), the mark \(m_{i}\), and the state of the LOB at the time of the event \(X_{t}\).
    \end{itemize}
\end{paragraph}

This combined approach allows to properly model such cases:
\begin{quote}
    A huge trade (large mark \(m_{i}\)) just executed. Because
    the bid-ask spread is currently wide (state \(X_{t}\)), this massive
    trade will trigger a cascade of aggressive cancel-and-replace orders
    (\(\phi_{X_{t_{i}}}\) becomes heavily excited). However, if the
    spread were narrow, that exact same size trade would trigger virtually
    no reaction.
\end{quote}
%------------------------------------------------------------------------------------
\section{Extension: The Neural Hawkes Process (NHP)}\label{sec:extension_neural_hawkes_process}
%------------------------------------------------------------------------------------
While computationally beautiful, implementing a State-Dependent
Marked Hawkes process presents a massive hurdle: Parameter Explosion.
If you track a limit order book with 10 distinct states (e.g.,
deciles of queue imbalance) and use 3 types of marks (small, medium,
and large orders), your kernel combinations multiply exponentially.
Estimating these parameters in real time using Maximum Likelihood
Estimation (MLE) or optimization methods like the Karush–Kuhn–Tucker
(KKT) conditions becomes computationally intensive for high-frequency
setups. To solve this, modern quantitative research frequently
utilizes Neural Hawkes Processes (NHP). Instead of computing explicit
mathematical formulas for every state-mark combination, deep learning
architectures (like Continuous-Time LSTMs) digest the raw sequence of
states and marks together, outputting a dynamically changing arrival
intensity natively.

%------------------------------------------------------------------------------------
\section{The Hawkes Process based framework for optimal market making}\label{sec:hawkes_process_framework}
%------------------------------------------------------------------------------------
Essentially the framework shall incorporate the accurate modelling of
events that impact the mid-price and the bid-ask spread, as these are
the two main reference points from which we desire to obtain the
optimal depths for the MM quoting policy. Further we assume that the
MM quoting policy shall not be continuous, as the optimization is
formulised as an optimal switching problem where the market maker
will be penalized every time he revises his quotes due to the
price-time priority of modern exchanges.\footnote{\cite[See][,
        p.8]{law2015pure}, this essentially assumes the MM can't reposition
    continuously.} For the switching regimes, we want to obtain the best
quote depths for the strategy.\footnote{\cite[see][,
        p.9]{law2015pure}} The effective arrival rate of the market orders
hitting the market maker is determined from the target market share
parameter $\rho$ chosen in advance by the market
maker.\footnote{\cite[see][, p.9]{law2015pure}} Volumes are used as
marks and the MM is allowed to liquidate excess inventory using
MO.\footnote{\cite[see][, p.9]{law2015pure}} The control problem is
now a combined (simultaneous) optimal switching and impulse control
problem under a pure-jump environment.\footnote{\cite[see][,
        p.9]{law2015pure}}

% State here the full model after having discussed all the extensions mathematically
% As the exact in detail specifications of the baseline intensity and the kernel functions
% are still open we will define them here in the subsections
% As the connection of the oder flow arrival modelling using the Hawkes Process
% to the mid price is crucial, we will also cover this here in more detail
%------------------------------------------------------------------------------------
\subsection{Time-inhomogeneous (Baseline) Intensities}\label{sub_sec:time_inhomogeneous_baseline_intensities}
%------------------------------------------------------------------------------------
As we have seen in the section
\hyperref[sub_sec:messages_and_cancellations]{3.2.6 Messages and
    Cancellations}, the intraday activity pattern is not constant, but
rather follows a certain pattern, with higher activity during certain
hours of the contract lifetime, which is also known as the \emph{Samuleson
effect}. This means that the baseline intensity of the Hawkes process,
which represents the underlying rate of order arrivals, has to be
dependent on the remaining contract lifetime, therefore being
non-stationary. By incorporating time-inhomogeneous baseline
intensities into our modeling framework, we can better capture the
dynamics of order arrivals and improve the performance of our market
making strategy. For the baseline intensity we assume:
\begin{equation}\label{eq:baseline_intensity}
    \mu_{j}(t) = \gamma_{j} e^{\eta_{j} t}, \quad j \in \{1, \ldots, 6\},
\end{equation}
where:\footnote{\cite[See ][, p.80]{luckner2021hawkes}}

Note the important practical caveat, that when the bid-ask spread is
equal to the contract’s tick size ($\delta$), it is not possible to
place LOs inside the bid-ask spread and the intensities of buy and
sell LOs placed into the bid-ask spread should be zero.
%------------------------------------------------------------------------------------
\subsection{Excitation Kernels}\label{sub_sec:excitation_kernels}
%------------------------------------------------------------------------------------
In the Hawkes-Poisson Process, the probability drop-off of an event
triggered by a past event is called ''kernel''.\footnote{Alternatively,
referred to as ''de-cay/exciting/fertility kernel'', \cite[see ][, p.19]{law2015pure} } Together with the
previously specified baseline intensities, the kernel functions
determine the overall intensity of the process. The kernel functions
can be generally of arbitrary functional shape, where the most common
choices are the exponential, power-law, and Rayleigh kernels. A
comparing figure of these kernel function types is provided below.

Here also a plot comparing the different popular choices for the
exciattion kernels, but also showing the spline based non-parametric
kernel estimation variant. After the empirical findings from section
XX, we chose to use the exponential kernel function, specified as:
\begin{equation}\label{eq:exponential_kernel}
    \phi_{j, k}(t) = \alpha_{j, k} e^{-\beta_{j, k}t}, \quad j, k \in \{1, \ldots, 6\},
\end{equation}
where:\footnote{\cite[See ][, p.80]{luckner2021hawkes}}
%------------------------------------------------------------------------------------
\subsection{Midprice Dynamics}\label{sub_sec:midprice_dynamics}
%------------------------------------------------------------------------------------
Having fully specified the order arrival rate dynamics, we can now turn
to the midprice dynamics. This part represents a crucial aspect of our
framework, as this often done using an arbitrary set process, that is
\emph{independent} of the order arrival events considered here, which is not
realistic. There are two possible scenarios, either there is enough
liquidity at both the best bid and the best ask levels or the
opposite (a just one sided scenario will be ignored). In the first
setting an incoming MO (supposedly even with bigger volume) will not
cause a change in the midprice, while in the second setting even a
small MO will probably cause either the bid or the ask side to be
depleted, causing the midprice to change. As we have seen in the
empirical section, the latter scenario is usually the default in the
CID market, as the liquidity tends to be low and the bid-ask spread several
tick sizes wide. Therefore we require a setting where the mid-price
is impacted by our mid-price moving event types and is of the type of
a pure jump process, following the considerations of \cite[see][,
    p.7]{law2015pure}. The directly resulting linkage between the best
bid and ask prices and incoming MO (that shall be connected) can be
seen here: \cite[see][, p.8]{law2015pure}

Therefore we assume it is best to not only submit orders to the best
bid and ask level, but rather consider a setting where the optimal
depths for posting orders should be determined, as the best bid and
ask levels can be depleted.

As stated previously, we concentrate on those events that move the
midprice. These are the \emph{aggressive} events that can be depicted
in Table XX and their empirical relative shares in Table XX. Namely
they are: placement of buy and sell LOs into the bid-ask spread,
depletion of the first or more levels of the LOB, and the placement
of buy and sell MOs. \footnote{\cite[See ][,
        p.80]{luckner2021hawkes}} Therefore resulting in the use of a
6-dimensional Hawkes process. As sanity check we have to note that in
the case where the bid-ask spread is just tick-size, it is of course
not possible to place further LO inside the spread, so the
probability of this event type is zero, and the model should be able
to capture this. The mid-price must be further connected to say at
least MO event types, for two reasons: First, otherwiese it can
happen that a sell MO arrives and the mid-price jumps into the
opposite direction if not coupled and secondly the MO carrying mark,
say its attached volume must be considered. A very big MO should have
a bigger impact on the mid-price as its probably gonna deplete the
best few levels until its filled than say a small MO, that
\textit{does not walk the book}. Further, we will extend the proposed
basis framework of \cite{luckner2021hawkes}, where his considered
main mid-price moving events are the following six: placement of buy
and sell LOs into the bid-ask spread, depletion of the first or more
levels of the LOB, and the placement of buy and sell MOs by the
following more granular extensions: bid-side MO, LO, cancellations at
L1, L2, L3. ask-side MO, LO, cancellations at L1, L2, L3.
% How should the baseline intensities be modelled
% Must be time varying to capture the intraday activity pattern, but also
% must be state dependent to capture the current state of the LOB, as
% the baseline intensity is the underlying rate of order arrivals, which
% is not only dependent on the time of the day, but also on the current state of the LOB, e.g. if the bid-ask spread is wide, the baseline
% How should the excitation kernel functions be chosen
% How can we integrate states with external variables
% How can we integrate marks with the order size and direction

Here the mathematical definition of the midprice dynamics, which is
the central variable of interest in our modeling framework. The

In reference to the before mentioned order type classification, only
informed traders post aggressive order types that move the mid price,
here from the data again measure the shares of these information
carrying order types and show if this problem is immanent and how we
can obtain the frequencies from these stats

We need to have the midprice impacted by incoming MO due to realism,
cause a certain jump of the process, jump size and direction have to
be realistic (in their direction, foresign and magnitude).

Must be a (pure) jump process since it is a discrete price grid, and
the midprice can only jump in discrete increments of the tick-size,
which is 0.01 EUR/MWh in the CID.

As the operating frequency at which LOs are posted inside the bid-ask
spread and therefore cause it and the mid-price to move is much lower
when compared with traditional financial markets and occurs randomly,
it is not uncommon to oberseve phases where there is no change at
all.\footnote{\cite[See ][, p.78]{luckner2021hawkes}} A natural
choice is therefore to model the mid-price and the bid-ask spread as
pure jump process.\footnote{\cite[See ][, p.78]{luckner2021hawkes}}

Based on \cite[See ][, p.262]{cartea2015algorithmic}, the midprice
dynamics can be used to account for adverse-selection risk a MM
faces, by either assuming that the midprice undergoes jumps in the
direction of incoming MO, or by letting the midprice's drift have a
short term alpha component which is affected by the arrival of MO.
%------------------------------------------------------------------------------------
\subsection{Cost functions}\label{sub_sec:cost_fuctions}
%------------------------------------------------------------------------------------


%------------------------------------------------------------------------------------
\section{Overview: Modelling Specifications}\label{sub_sec:overview_modelling_specs}
%------------------------------------------------------------------------------------
Although possible in theory, the modelling of cross-excitement among
various hourly products is not pursued in this thesis, as the
in-parallel lifetime of the contracts differs (See Table) what makes
it hard to calculate the cross-excitation, and [Zweiter Grund].

Ignoring the external influence of updated forecasts and other
external variables, can be partly defended by seeing that fundamental
drivers are partly not relevenat for the liquidity, see here:
\cite{hagemann2013empirical}.

Also the idea, i.e when we refit daily, we could use these values as
prior for the next fit making use of Bayesian updating, see
\cite{jain2024limit}

Here a table as overview of all different model versions that we fit
and want to compare their performances in judging whether some
specifications and inclusions are justified or not.

Having discussed the rich possibilities in that the modelling
frameworkmcan be specified and extended, we will now turn to the
specifications of the modelling frameworks we want to test:
\begin{enumerate}
    \item We start off by the model proposed by \cite{luckner2021hawkes}, which
          is a 6-dimensional Hawkes process, where the midprice moving events
          are the following six: placement of buy and sell LOs into the bid-ask
          spread, depletion of the first or more levels of the LOB, and the
          placement of buy and sell MOs.
    \item Next we separate the bid/ask event directions, considering: bid-side
          aggressive events, ask-side aggressive events, side-specific arrival
          intensities and side-specific impact on midprice and spread. This
          should serve us to test whether asymmetry between buy and sell flow
          matters.
    \item On top the event types are added: market orders, limit orders,
          cancellations
    \item Additionally we will consider not only the best bid/ask level but
          rather: L1, L2, L3 on bid side, L1, L2, L3 on ask side,
          level-specific intensities and excitation and level-specific queue
          depletion and replenishment.
    \item Letting intensities react to the current book state is achieved by
          adding state dependencies using possible variables like: spread,
          midprice, imbalance, queue sizes at L1–L3, time-of-day /
          time-to-delivery, possibly volatility regime.
    \item We will next add marks to incorporate event magnitudes with the
          following possible components: order size, price impact size,
          possibly depth consumed, possibly aggressiveness of the order.
    \item In a more optional setting we can benchmark the so far explored
          settings against a Reinforcement Learning extension, with this broad
          formulation: Hawkes simulator, state = intensities + book state +
          inventory, actions = bid/ask quote placement, reward = PnL minus risk
          penalties.
\end{enumerate}

The specifications in that order look like follows. The starting base
model of \cite{luckner2021hawkes} is a 6-dimensional Hawkes process.
The vector of event counting processes $N(t)$ is given as:
\begin{equation}\label{eq:vector_event_counting_processes}
    N(t) = (N^{BM}(t)m, N^{SM}(t), N^{BL, +}(t), N^{BL, -}(t), N^{SL, +}(t), N^{SL, -}(t))
\end{equation}
where
\begin{itemize}
    \item $N^{BM}(t)$ is the counting process for buy market orders,
    \item $N^{SM}(t)$ is the counting process for sell market orders,
    \item $N^{BL, +}(t)$ is the counting process for buy limit orders placed inside the bid-ask spread,
    \item $N^{BL, -}(t)$ is the counting process for buy limit orders placed outside the bid-ask spread,
    \item $N^{SL, +}(t)$ is the counting process for sell limit orders placed inside the bid-ask spread, and
    \item $N^{SL, -}(t)$ is the counting process for sell limit orders placed outside the bid-ask spread.
\end{itemize}
For each of the six processes: $i \in \{1, \ldots, 6\}$ define the conditional intensities:
\begin{equation}\label{eq:conditional_intensities}
    \lambda_{i}(t) = \mu_{i}(t, X_{t}) + \sum_{j=1}^{6} \int_{0}^{t} \phi_{ij}(t - s) dN_{j}(s)
\end{equation}
where:
\begin{itemize}
    \item $\lambda_{i}(t)$ conditional arrival intensity of event type $i$,
    \item $\mu_{i}(t, X_{t})$ is the baseline intensity for process $i$, which can depend on time $t$ and the current state of the limit order book $X_{t}$,
    \item $\phi_{ij}(t - s)$ is the excitation kernel function that captures the influence of from event type $j$ to event type $i$,
    \item $X_{t}$ state vector, e.g. spread, midprice, imbalance, time-to-close at time $t$.
\end{itemize}
The mid-price and spread dynamics can be formulated as follows.
Define the mid-price as $P_{t}$ and the Spread as $S_{t}$ as pure jump processes, driven
by the price-moving events:
\begin{equation}
    dP_{t} = \sum_{i=1}^{6} \Delta_{i}^{P} dN_{i}(t)
\end{equation}
and the spread as:
\begin{equation}
    dS_{t} = \sum_{i=1}^{6} \Delta_{i}^{S} dN_{i}(t)
\end{equation}
where
\begin{itemize}
    \item $\Delta_{i}^{P}$ is the price impact of event type $i$ on the mid-price,
    \item $\Delta_{i}^{S}$ impact of event $i$ on the spread.
\end{itemize}

State dependency can be added by describing the state using the
possible variables combined as:
\begin{equation}
    X_{t} = (S_{t}, P_{t}, I_{t}, Q_{t}^{L1}, Q_{t}^{L2}, Q_{t}^{L3}, T_{t}, V_{t})
\end{equation}
where:
\begin{itemize}
    \item $S_{t}$ is the bid-ask spread at time $t$
    \item $P_{t}$ is the mid-price at time $t$
    \item $I_{t}$ is the order imbalance at time $t$
    \item $Q_{t}^{L1}, Q_{t}^{L2}, Q_{t}^{L3}$ are the queue sizes at levels 1, 2, and 3 of the limit order book at time $t$
    \item $T_{t}$ is the time-to-delivery at time $t$
    \item $V_{t}$ is the volatility regime at time $t$
\end{itemize}

Then a baseline can be chosen as:
\begin{equation}
    \mu_{i}(t, X_{t}) = f_{i}(S_{t}, P_{t}, I_{t}, Q_{t}^{L1}, Q_{t}^{L2}, Q_{t}^{L3}, T_{t}, V_{t})
\end{equation}
or alternatively:
\begin{equation}
    \mu_{i}(t, X_{t}) = exp(\alpha_{i} + \beta_{i}^{\top}X_{t}+g_{i}(t))
\end{equation}
where:
\begin{itemize}
    \item $\alpha_{i}$ channel-specific level,
    \item $\beta_{i}$ sensitivity to the state variables,
    \item $g_{i}(t)$ deterministic intraday seasonality..
\end{itemize}

Including marks $m_{k}$ to each event, for example the order size as
mark, a marked kernel can be written as:
\begin{equation}
    \phi_{ij}(u, m) = a_{i, j} h(m) e^{-\omega_{i, j} u} \mathbb{1}_{u > 0}
\end{equation}
where:
\begin{itemize}
    \item $a_{i, j}$ strength of excitation,
    \item $\omega_{i, j}$ decay speed,
    \item $h(m)$ mark impact function, e.g. $h(m) = m$ or $h(m) = \log(1+m)$
\end{itemize}
what lets larger MOs or LOs have bigger effects than smaller ones.

The level resolved extension results from refining each aggregated
event type further by side, event type and level, resulting in a
18-dimensional Hawkes process: $i \in$ \{side x type x level\} with:
\begin{itemize}
    \item side $\in$ \{bid, ask\},
    \item type $\in$ \{MO, LO, CA\},
    \item level $\in$ \{1, 2, 3\}.
\end{itemize}
with the intensities becoming:
\begin{equation}
    \lambda_{i}(t) = \mu_{i}(t, X_{t}) + \sum_{j=1}^{18} \int_{0}^{t} \phi_{ij}(t - s, m) dN_{j}(s),\\
    i = 1, \ldots, 18
\end{equation}

Stability conditions: For exponential kernels, define the integrated
excitation matrix $G$ as:
\begin{equation}
    G_{ij} = \int_{0}^{\infty} \phi_{ij}(u) du = \frac{a_{i, j}}{\omega_{i, j}}
\end{equation}
A sufficient stability condition is $\rho(G) < 1$
with:
\begin{itemize}
    \item $\rho(G)$ is the spectral radius of the matrix $G$,
    \item $G_{ij}$ is the expected number of events of type $i$ triggered by a single event of type $j$.
\end{itemize}

%------------------------------------------------------------------------------------
\section{Deriving the optimal strategy: the optimal depths}\label{sec:optimal_depths}
%------------------------------------------------------------------------------------
Here the derivation and solution of the stochastic optimal control
problem for the optimal depths, which is

%------------------------------------------------------------------------------------
\subsection{Stochastic Optimal Control}\label{sub_sec:stochastic_optimal_control}
%------------------------------------------------------------------------------------
The goal of stochastic control problems is to maximise (minimise)
some expected profit (cost) function by tuning a strategy which
itself affects the dynamics of the underlying stochastic system and
to find the strategy which attains the maximum
(minimum).\footnote{\cite[See ][, p.100]{cartea2015algorithmic}}

A prominent tool for solving such problems is the \emph{dynamic
    programming principle} (DPP) and the related non-linear partial
differential equation (PDE), also knwon as
\emph{Hamilton-Jacobi-Bellman} (HJB) equation.\footnote{\cite[See ][,
        p.100]{cartea2015algorithmic}} With the DPP a stochastic control
problem can be solved by breaking it down into a sequence of smaller
subproblems, which can be solved recursively. The HJB equation
provides a necessary condition for optimality and can be used to
derive the optimal control strategy by solving the PDE subject to
appropriate boundary conditions.\footnote{\cite[See ][,
        p.100]{cartea2015algorithmic}}

In the context of this thesis, we are especially interested in the
setting where a MM wishes to maximise expected terminal wealth by
using MO and LO.
%------------------------------------------------------------------------------------
\subsubsection{Control for Diffusion Processes}\label{sub_sub_sec:control_for_diffusion_processes}
%------------------------------------------------------------------------------------

%------------------------------------------------------------------------------------
\subsubsection{Control for Counting Processes}\label{sub_sub_sec:control_for_counting_processes}
%------------------------------------------------------------------------------------

%------------------------------------------------------------------------------------
\section{Estimation techniques}\label{sec:estimation_techniques}
%------------------------------------------------------------------------------------
Here describe the maximum likelihood or GMM estimations for the
various model parameters, before we head over to the next pratical
chapter where we apply all together.
%------------------------------------------------------------------------------------
\chapter{Numerical Analysis and Simulation}\label{ch:numerical_analysis}
%------------------------------------------------------------------------------------
Here the concrete estimation procedures for the arrival rates and
intensities
%------------------------------------------------------------------------------------
\section{Modelling Order arrivals}\label{sec:modelling_order_arrivals}
%------------------------------------------------------------------------------------
Since the LOB and every state can essentially be reconstructed by the
two main input parameters, that are the arrival of limit orders and
market orders, we can focus on the modelling of these two main order
types, as well as cancellations, which are the main drivers of the
LOB dynamics. From these general order types, we further focus on
those that are actually relevant, i.e. those that move the mid price,
which are the so called \textit{aggressive orders}, according to the
detailed order-categorization based on XYZ laid out earlier in
chapter XYZ. Another central and shaping input is the volume depth of
the LOB.\cite{choi2021optimal} The basic idea behind modelling order
arrivals is the following one: some event types are stochastically
excited or inhibited—in the sense of having their probabilities
elevated or decreased—by patterns in the sequence of previous
events.\footnote{\cite[See ][, p.1]{mei2017neural}} Learning this
pattern enables us to predict which type of event will happen next
and when.\footnote{\cite[See ][, p.1]{mei2017neural}}

Remembering the critical consideration, that a MM has no intrinsic
need or desire to hold any inventory, so she will only buy (sell) in
anticipation of a subsequent sale (purchase), what motivates the
pursued modelling of order arrivals.\footnote{\cite[See ][,
        p.21]{cartea2015algorithmic}}

The estimated relative shares per order type here from the data:
%------------------------------------------------------------------------------------
%	                  FIGURE - START
%------------------------------------------------------------------------------------
\begin{figure}[htbp]
    \vspace*{0.5cm}
    \centering
    \blendfigure{1.0\textwidth}{7cm}{rel_share_order_types.png}
    \captionfootcite{Relative Percentage share of each order type.}{Based on \cite{law2015pure}, own illustration.}\label{fig:rel_share_order_types}
\end{figure}
%------------------------------------------------------------------------------------
%	                  FIGURE - END
%------------------------------------------------------------------------------------

%------------------------------------------------------------------------------------
\subsection{Market Order Arrivals}\label{sub_sec:mo_arrivals}
%------------------------------------------------------------------------------------
\subsection{Limit Order Arrivals}\label{sub_sec:lo_arrivals}
%------------------------------------------------------------------------------------
\subsection{Cancelations}\label{sub_sec:cancellations}
%------------------------------------------------------------------------------------
%------------------------------------------------------------------------------------
\subsection{The role of Order Imbalance}\label{sub_sec:order_imbalance}
%------------------------------------------------------------------------------------
%------------------------------------------------------------------------------------
\section{Exchange Latency}\label{sec:exchange_latency}
%------------------------------------------------------------------------------------
Here the exact calculation of the exchange latency is provided, which
is the time it takes for an order to be processed and executed on the
exchange after it has been submitted. This includes the time taken
for the order to reach the exchange, be processed by the exchange's
matching engine, and receive a confirmation of execution. The
exchange latency can vary based on several factors, including the
type of order, the current market conditions, and the specific
exchange being used. Understanding and minimizing exchange latency is
crucial for developing effective trading strategies, particularly in
high-frequency trading environments where even small delays can
impact profitability. important for sampling and fitting durations
etc.

%------------------------------------------------------------------------------------
\section{Goodness of Fit Evaluation}\label{sec:goodness_of_fit_evaluation}
%------------------------------------------------------------------------------------
QQ Plots here, AIC table, etc.
%------------------------------------------------------------------------------------
\section{Optimal Strategy Backtests}\label{sec:optimal_strategy_backtests}
%------------------------------------------------------------------------------------
Here some example graph that is shown in the List of Figures:
\begin{figure}[htbp]
    \centering
    \includegraphics[width=0.2\textwidth]{siegel_maths.png}
    \captionfootcite{Example Figure.}{Test cite \cite{law2019market}, as of 16.05.2026.}\label{fig:example_figure}
\end{figure}

\begin{figure}[htbp]
    \centering
    \blendfigure{8cm}{8cm}{siegel_maths.png}
    % \caption{Example Figure with blended background}\label{fig:example_figure_blended}
    \captionfootcite{Some text.}{See first: \cite{eeg_2023_feed_in}.}
\end{figure}

Here some example table that will be shown in the List of Tables
section:
\begin{table}[htbp]
    \centering
    \begin{tabular}{|c|c|c|}
        \hline
        Column 1 & Column 2 & Column 3 \\
        \hline
        Data 1   & Data 2   & Data 3   \\
        Data 4   & Data 5   & Data 6   \\
        \hline
    \end{tabular}
    \caption{Example Table}\label{tab:example_table}
\end{table} listing example that will be shown in the List of Algorithms section:
\begin{lstlisting}[language=Python, caption={Example Algorithm}]
def example_algorithm(data):
    # Some example algorithm code
    result = data * 2
    return result
\end{lstlisting}

Her some example algorithm that will appear in the List of Algorithms
section:
\begin{algorithm}[htbp]
    \caption{Example Algorithm}
    \label{alg:example_algorithm}
    \begin{algorithmic}[1]
        \Require $data$ - Input data for the algorithm
        \Ensure $result$ - Output result of the algorithm
        \State Initialize $result$ to 0
        \For{each element in $data$}
        \State Update $result$ based on the element
        \EndFor
        \State \Return $result$
    \end{algorithmic}
\end{algorithm}

Here some examples of different citation styles offered by \LaTeX:
\textcite{antonio2022inflation}. Parenthetical citation:
\parencite{antonio2022inflation}. Footnote citation:
\autocite{antonio2022inflation}. Footnote citation with page number:
\autocite[1]{antonio2022inflation}. Inline citation:
\textcite{rosenberg_2003_exchange_rate_determination}. Inline
citation: \textcite{antonio2022inflation}. Inline citation:
\textcite{seabold2010statsmodels}. Inline citation:
\textcite{seabold2010statsmodels}. Inline citation:
\textcite{fama2016commodity}. Inline citation:
\textcite{rosenberg_2003_exchange_rate_determination}. Inline
citation: \textcite{fama2016commodity}. Example of a footnote
citation.\footnote{\cite{fama2016commodity}}.

An example of an equation (numbered and centered) is shown below:
\begin{equation}
    E = mc^2
\end{equation}
%------------------------------------------------------------------------------------
%	                  LITERATURE REVIEW
%------------------------------------------------------------------------------------
% \chapter{Literature Review}
% \label{ch:literature_review}
% % 02-literature-review
% % Here the introduction text
% Here the Literature review text.
% \section{Example Section}
%------------------------------------------------------------------------------------
%	                  METHODOLOGY
%------------------------------------------------------------------------------------
% \chapter{Methodology}
% \label{ch:methodology}
% % 03-methodology
% % Here the introduction text
% Here the methodology text.
%------------------------------------------------------------------------------------
%	                  DATA
%------------------------------------------------------------------------------------
% \chapter{Data}
% \label{ch:data}
% % 04-data
% % Here the introduction text
% Here the data text.
%------------------------------------------------------------------------------------
%	                  EXPERIMENTS
%------------------------------------------------------------------------------------
% \chapter{Experiments}
% \label{ch:experiments}
% % 05-experiments
% % Here the introduction text
% Here the experiments text.
%------------------------------------------------------------------------------------
%	                  RESULTS
%------------------------------------------------------------------------------------
\chapter{Results}
\label{ch:results}
% 06-results
% Here the introduction text
Here the results text.
%------------------------------------------------------------------------------------
%	                  DISCUSSION
%------------------------------------------------------------------------------------
\chapter{Discussion}
\label{ch:discussion}
% 07-discussion
% Here the introduction text
Here the discussion text.
%------------------------------------------------------------------------------------
%	                  CONCLUSION
%------------------------------------------------------------------------------------
\chapter{Conclusion}
\label{ch:conclusion}
% 08-conclusion
% Here the introduction text
Here the conclusion text.
%------------------------------------------------------------------------------------
%	                  REFERENCES
%------------------------------------------------------------------------------------
\chapter*{Bibliography}
\markboth{Bibliography}{}
\addcontentsline{toc}{chapter}{Bibliography}
% \printbibliography[heading=none]
\printbibliography[type=article,heading=subbibliography,title={Articles}]
\printbibliography[type=book,heading=subbibliography,title={Books}]
\printbibliography[type=report,heading=subbibliography,title={Reports and Working Papers}] % for techreports
\printbibliography[keyword={online},heading=subbibliography,title={Online Sources}]
\printbibliography[keyword={data},heading=subbibliography,title={Data Sources}]
\printbibliography[type=inproceedings,heading=subbibliography,title={Conference Proceedings}]
\printbibliography[type=incollection,heading=subbibliography,title={Book Chapters}]
\printbibliography[type=unpublished,heading=subbibliography,title={Lecture notes}]
\printbibliography[keyword={litreview},heading=subbibliography,title={Full sources for the literature overview}]
\printbibliography[type=misc,heading=subbibliography,title={Reports and Online Sources}]
%------------------------------------------------------------------------------------
%	                  APPENDIX
%------------------------------------------------------------------------------------
\begin{appendix}\label{anhang}
    \chapter{Supplementary Market Structure Information}\label{anhang1}
    \section{Hourly Products and Trading Durations}\label{sec:hourly_products_and_trading_durations}
    \begin{table}[htbp]
        \centering
        \tiny
        \resizebox{\textwidth}{!}{%
        \begin{tabular}{l c c c c c c c}
            \toprule
            \textbf{Hourly}  & \textbf{Delivery}       & \textbf{Cross-Border} & \textbf{Max Cross-Border} & \textbf{Local German} & \textbf{Max Local German} & \textbf{Next-Day Hours}   & \textbf{Duration of This} \\
            \textbf{Product} & \textbf{Interval (CET)} & \textbf{Gate Closure} & \textbf{Trading Duration} & \textbf{Gate Closure} & \textbf{Trading Duration} & \textbf{Open in Parallel} & \textbf{Parallel Window}  \\
            \midrule
            Hour 01          & 00:00 -- 01:00          & 23:30 (D-1)           & 8 h 30 min                & 23:55 (D-1)           & 8 h 55 min                & 24                        & 8 h 55 min                \\
            Hour 02          & 01:00 -- 02:00          & 00:30 (D)             & 9 h 30 min                & 00:55 (D)             & 9 h 55 min                & 23                        & 1 h 00 min                \\
            Hour 03          & 02:00 -- 03:00          & 01:30 (D)             & 10 h 30 min               & 01:55 (D)             & 10 h 55 min               & 22                        & 1 h 00 min                \\
            Hour 04          & 03:00 -- 04:00          & 02:30 (D)             & 11 h 30 min               & 02:55 (D)             & 11 h 55 min               & 21                        & 1 h 00 min                \\
            Hour 05          & 04:00 -- 05:00          & 03:30 (D)             & 12 h 30 min               & 03:55 (D)             & 12 h 55 min               & 20                        & 1 h 00 min                \\
            Hour 06          & 05:00 -- 06:00          & 04:30 (D)             & 13 h 30 min               & 04:55 (D)             & 13 h 55 min               & 19                        & 1 h 00 min                \\
            Hour 07          & 06:00 -- 07:00          & 05:30 (D)             & 14 h 30 min               & 05:55 (D)             & 14 h 55 min               & 18                        & 1 h 00 min                \\
            Hour 08          & 07:00 -- 08:00          & 06:30 (D)             & 15 h 30 min               & 06:55 (D)             & 15 h 55 min               & 17                        & 1 h 00 min                \\
            Hour 09          & 08:00 -- 09:00          & 07:30 (D)             & 16 h 30 min               & 07:55 (D)             & 16 h 55 min               & 16                        & 1 h 00 min                \\
            Hour 10          & 09:00 -- 10:00          & 08:30 (D)             & 17 h 30 min               & 08:55 (D)             & 17 h 55 min               & 15                        & 1 h 00 min                \\
            Hour 11          & 10:00 -- 11:00          & 09:30 (D)             & 18 h 30 min               & 09:55 (D)             & 18 h 55 min               & 14                        & 1 h 00 min                \\
            Hour 12          & 11:00 -- 12:00          & 10:30 (D)             & 19 h 30 min               & 10:55 (D)             & 19 h 55 min               & 13                        & 1 h 00 min                \\
            Hour 13          & 12:00 -- 13:00          & 11:30 (D)             & 20 h 30 min               & 11:55 (D)             & 20 h 55 min               & 12                        & 1 h 00 min                \\
            Hour 14          & 13:00 -- 14:00          & 12:30 (D)             & 21 h 30 min               & 12:55 (D)             & 21 h 55 min               & 11                        & 1 h 00 min                \\
            Hour 15          & 14:00 -- 15:00          & 13:30 (D)             & 22 h 30 min               & 13:55 (D)             & 22 h 55 min               & 10                        & 1 h 00 min                \\
            Hour 16          & 15:00 -- 16:00          & 14:30 (D)             & 23 h 30 min               & 14:55 (D)             & 23 h 55 min               & 9                         & 1 h 00 min                \\
            Hour 17          & 16:00 -- 17:00          & 15:30 (D)             & 24 h 30 min               & 15:55 (D)             & 24 h 55 min               & 8                         & 1 h 00 min                \\
            Hour 22          & 21:00 -- 22:00          & 20:30 (D)             & 29 h 30 min               & 20:55 (D)             & 29 h 55 min               & 3                         & 1 h 00 min                \\
            Hour 23          & 22:00 -- 23:00          & 21:30 (D)             & 30 h 30 min               & 21:55 (D)             & 30 h 55 min               & 2                         & 1 h 00 min                \\
            Hour 24          & 23:00 -- 24:00          & 22:30 (D)             & 31 h 30 min               & 22:55 (D)             & 31 h 55 min               & 1                         & 1 h 00 min                \\
            \bottomrule
        \end{tabular}}
        \vfill
        \footnotesize\textit{Note: All 24 next-day hours open together at 15:00 CET (D-1). The count decreases as each individual product hits its respective local gate closure.}
        \captionfootcite{EPEX Spot CID Trading Durations and Parallel Product Counts.}{Based on \cite{kiesel_efc19_2019}, as of 16.05.2026.}\label{tab:epex_parallel_durations}
    \end{table}

    \chapter{Data}\label{anhang2}
    \section{LOB Data Description}\label{sec:lob_data_description}
    \begin{table}[htbp]
        \centering
        \resizebox{\textwidth}{!}{%
        \begin{tabular}{|c|c|}
            \hline
            Column Name       & Description                                                                    \\
            \hline
            Timestamp         & The date and time of the order event                                           \\
            Order Type        & The type of order (e.g., limit order, market order, cancellation)              \\
            Price             & The price at which the order was placed or executed                            \\
            Quantity          & The quantity of the order                                                      \\
            Side              & The side of the order (buy or sell)                                            \\
            Order ID          & A unique identifier for the order                                              \\
            Trader ID         & A unique identifier for the trader who placed the order                        \\
            Venue             & The trading venue where the order was placed                                   \\
            Status            & The status of the order (e.g., filled, partially filled, canceled)             \\
            Time in Force     & The duration for which the order is valid (e.g., GTC, IOC)                     \\
            Execution ID      & A unique identifier for the execution event (if applicable)                    \\
            Counterparty ID   & A unique identifier for the counterparty involved in the trade (if applicable) \\
            Order Book State  & A snapshot of the limit order book at the time of the event (optional)         \\
            Market Conditions & Additional market conditions at the time of the event (optional)               \\
            Notes             & Any additional notes or comments related to the order event (optional)         \\
            \hline
        \end{tabular}}
        \captionfootcite{Data Columns and Their Descriptions.}{Based on the typical structure of order book data, as of 16.05.2026.}\label{tab:data_columns}
    \end{table}
    % Here the first part of the appendix text
    Here the first part of the appendix text.
    \section{Second Part}
    % Here the second part of the appendix text
    Here the second part of the appendix text.

    \chapter{Appendix B}\label{anhang3}

    \chapter{Appendix C}\label{anhang4}

    \chapter{Code}\label{anhang5}
    Here some coding snippets:
    \begin{lstlisting}[language=Python, caption={Example Python Code}]
    def example_function(x):
        return x**2
    \end{lstlisting}

    \chapter{Engineering Configurations}\label{anhang6}
    \section{Hyperparameter Optimization}

    \section{Computation Time}

\end{appendix}
%------------------------------------------------------------------------------------
%	                  AFFIRMATION
%------------------------------------------------------------------------------------
\chapter*{Affirmation} % by using * we are not showing this chapter in the table of contents
\markboth{~}{} % Clear header for the table of contents
\addcontentsline{toc}{chapter}{Affirmation}
\textbf{General Declaration of Authorship}\\
I hereby declare that I have composed my Master’s thesis
\makeatletter\textit{``\@title''}\makeatother\ independently using
only those resources mentioned, and that I have as such identified
all passages which I have taken from publications verbatim or in
substance. I agree that the work will be reviewed using plagiarism
testing software. Neither this paper, nor any extract of it, has been
previously submitted to an examining authority, in this or a similar
form. I have ensured that the written version of this thesis is
identical to the version saved on the enclosed storage medium.

\textbf{Use of Artificial Intelligence}
\begin{itemize}
    \item[a)] If the use of generative artificial intelligence systems as an
              aid is prohibited, I declare that I have not used content
              generated by such systems in preparing this thesis.
    \item[b)] If the use of generative artificial intelligence systems as an
              aid is permitted, I declare that I have used such systems only
              as explicitly authorised by the lecturer or examiner and have
              fully documented all generated outputs used in this thesis.
\end{itemize}

\begin{flushright}
    \vspace{2cm} % Adjust vertical spacing as needed
    \rule{10cm}{0.4pt} \\ % Horizontal line for the signature
    Place, Date, Signature \\[1cm] % Label for the signature line
\end{flushright}
\end{document}